\paragraph{About this part.} Part~II surveys the models of the order book, from the simplest to the
most complex. It starts with the empirical facts any model must respect (Section~\ref{sec:facts});
moves through hand-built signals computed from the book (order-flow imbalance, queue imbalance, the
micro-price; Section~\ref{sec:ofi}); then to models of the queues themselves and the time until they
empty (Section~\ref{sec:queue-models}); to Hawkes processes, which model how events trigger further
events (Section~\ref{sec:hawkes}); to deep networks (Section~\ref{sec:deep}); to hybrids of the two
(Section~\ref{sec:hybrid}); and finally to information from related instruments
(Section~\ref{sec:cross}). Every model ends with a model card answering the same four questions.

% ================================================================
\section{Stylised facts a model should reproduce}
\label{sec:facts}
% ================================================================

Before turning to models, it helps to list what any adequate model of the book must be consistent
with. These regularities have been documented repeatedly across markets and periods.
\begin{description}[leftmargin=1.5em,style=nextline,itemsep=2pt]
\item[Events cluster] Order-book events arrive in bursts separated by quiet periods, and a large share
  of events are triggered by earlier ones (Section~\ref{sec:hawkes}).
\item[Long waits are common] Times between price changes have heavy tails; in a simple model of exactly
  balanced queues the expected waiting time is even infinite~\citep{contlarrard2013}.
\item[Order flow has long memory] The signs of successive market orders are positively correlated
  over very long lags, with a Hurst exponent of about $0.7$, mainly because large orders are split into
  many small ones~\citep{lillo2004,lillo2005}.
\item[Prices nevertheless diffuse] Despite persistent order flow, prices are close to unpredictable
  beyond short horizons, because limit orders absorb the pressure: the market is held near a critical
  balance between the two~\citep{bouchaud2004}.
\item[Short-horizon mean reversion] Measured volatility is higher over very short windows than over
  longer ones (bid--ask bounce), and correlations between assets fall when measured over very short
  windows (the Epps effect)~\citep{bacry2013}.
\item[The average book is humped] Average depth is largest a few ticks away from the best price, not
  at it~\citep{bouchaud2002}.
\item[Impact is concave] The price impact of a large order executed in pieces grows less than linearly
  with its size, close to a square-root law~\citep{bouchaud2018}; Hawkes models reproduce its concave, relaxing
  shape~\citep{bacry2014}.
\item[Order flow explains price changes] Over short windows, price changes are approximately linear in
  order-flow imbalance, with a slope inversely related to depth~\citep{cks2014}.
\item[Not only trades move prices] Roughly half of price-changing events are limit orders placed inside
  the spread~\citep{eisler2012}, and the book beyond the best quotes carries part of the price
  information~\citep{cao2009}.
\item[Activity is seasonal] Intensities and volatility follow intraday patterns, so models fitted with a
  constant baseline over a non-stationary period overstate how self-exciting the market is~\citep{filimonovsornette2015}.
\end{description}
Generative models are judged by how many of these they reproduce (Section~\ref{sec:deep-gen});
predictive models are judged by whether they exploit them after costs.

\paragraph{How each model is presented.} Every model in this part ends with a \emph{model card} that
answers the same four questions: what has to be calibrated (the numbers estimated from data), how it is
calibrated, what the fitted model is then used for, and a picture that makes it easy to remember and
to teach. The cards summarise the section they close; the detail and the citations are in the text
above each one.

% ================================================================
\section{Order-flow imbalance and queue imbalance}
\label{sec:ofi}
% ================================================================

The simplest useful models of the book are single numbers computed directly from it. This section
covers the two most important: \emph{order-flow imbalance}, which measures the net buying pressure
from the stream of book events, and \emph{queue imbalance}, which measures how lopsided the two best
queues are at an instant. It then refines queue imbalance into better estimates of the ``true'' price
(the weighted mid and the micro-price) and describes features of the book's shape beyond the best
level. These signals are cheap, interpretable and hard to beat, which makes them the baselines for
everything that follows.

\subsection{Order-flow imbalance}

\paragraph{The idea.} Why does the price move? Look again at Figure~\ref{fig:book}. The best ask
is 5000.25 with 14 lots waiting. If buyers take all 14 lots, or the sellers waiting there cancel,
the best ask becomes 5000.50 and the mid rises. If instead more sellers join at 5000.25, the queue
grows and a rise becomes harder. The same holds in mirror image on the bid side. So anything that
\emph{grows the bid queue or shrinks the ask queue} pushes the price up, and anything that
\emph{shrinks the bid queue or grows the ask queue} pushes it down. Order-flow imbalance (OFI)
simply keeps a running score of these pushes~\citep{cks2014}.

\paragraph{Measuring one event.} Visible changes at the best prices come from three event types:
limit orders arriving, aggressive orders executing against waiting ones, and cancellations. For
event number $n$, define
\[
  \Delta v^{b}_{n} = v^{b}_{1,t_n} - v^{b}_{1,t_n-},
  \qquad
  \Delta v^{a}_{n} = v^{a}_{1,t_n} - v^{a}_{1,t_n-}.
\]
Recall that $v^{b}_{1,t}$ is the number of lots waiting at the best bid at time $t$, and that
$t_n-$ means ``just before event $n$''. So $\Delta v^{b}_{n}$ is \emph{the size of the best-bid
queue just after event $n$ minus its size just before}: how many lots event $n$ added to the
best-bid queue (positive) or removed from it (negative). $\Delta v^{a}_{n}$ is the same for the
best-ask queue. In the book of Figure~\ref{fig:book}:
\begin{itemize}[leftmargin=1.5em,itemsep=1pt]
\item a new 20-lot buy order at 5000.00 gives $\Delta v^b_n = +20$ (29 lots become 49), and
  $\Delta v^a_n = 0$;
\item a sell market order for 10 lots trades against the front of the bid queue and gives
  $\Delta v^b_n = -10$;
\item a seller cancelling 3 lots at 5000.25 gives $\Delta v^a_n = -3$ (14 lots become 11).
\end{itemize}

\paragraph{Scoring and adding up.} While the best prices do not change, the score of event $n$
and the order-flow imbalance over the last $w$ seconds are
\begin{equation}
  e_n = \Delta v^{b}_{n} - \Delta v^{a}_{n},
  \qquad
  \OFI^{(w)}_t = \sum_{n:\, t_n \in (t-w,\, t]} e_n .
\end{equation}
\inwords{$e_n$ is positive when event $n$ grew the bid queue or shrank the ask queue---a push up---and
negative for a push down; subtracting $\Delta v^{a}_{n}$ is what turns a shrinking ask queue into a
positive score. The sum runs over every event $n$ whose time $t_n$ falls in the last $w$ seconds,
so $\OFI^{(w)}_t$ is the net buying pressure over that window, like the net flow into a tank: lots
in on the buy side and out on the sell side count as inflow, the reverse as outflow.}

When a best price \emph{does} change, the full definition of~\citet{cks2014} counts whole queues.
If the best bid rises to a new, higher price, the entire new queue counts as added buying interest;
if the best-bid level is emptied and the best bid falls, the entire old queue counts as removed. The
ask side is treated symmetrically. This keeps the score meaningful across price changes.

\paragraph{From pressure to price.} The associated short-horizon relation is a straight line,
\begin{equation}
  m_t - m_{t-w} = \beta\, \OFI^{(w)}_t + \varepsilon_t,
\end{equation}
where $\beta$ is the price-impact coefficient and $\varepsilon_t$ a residual.
\inwords{over a short window, the mid moves roughly in proportion to the net buying pressure.
$\beta$ converts lots of pressure into price: if $\beta$ is $0.01$ points per lot, a net 100 lots of
pressure moves the mid by about one point. $\varepsilon_t$ is whatever the line does not explain.
Fitting $\beta$ is an ordinary linear regression of price change on OFI.}
\citet{cks2014} find that $\beta$ is inversely related to market depth: in a thin book a few lots
empty a queue and move the price, while in a deep book the same lots barely dent it. OFI is
low-dimensional, cheap to compute and easy to interpret, which makes it the natural baseline
against which more complex models should be judged.

\begin{example}[OFI by hand]
\label{ex:ofi}
Five events at unchanged best prices, all inside one window of length $w$:
\begin{center}\small
\begin{tabular}{clrrr}
\toprule
$n$ & event & $\Delta v^b_n$ & $\Delta v^a_n$ & $e_n$ \\
\midrule
1 & bid limit order, 20 lots & $+20$ & $0$ & $+20$ \\
2 & ask cancel, 15 lots & $0$ & $-15$ & $+15$ \\
3 & buy market order hits ask, 10 lots & $0$ & $-10$ & $+10$ \\
4 & ask limit order, 30 lots & $0$ & $+30$ & $-30$ \\
5 & bid cancel, 5 lots & $-5$ & $0$ & $-5$ \\
\midrule
 & & & $\OFI^{(w)}_t$ & $+10$ \\
\bottomrule
\end{tabular}
\end{center}
Events 2 and 3 enter through identical arithmetic, yet one is a cancellation and the other an
aggressive trade. OFI does not distinguish them; Section~\ref{sec:events} returns to this.
\end{example}

\subsection{Evidence and extensions}

\paragraph{How well OFI explains price changes.} On 50 US stocks, \citet{cks2014} find that the
linear relation explains mid-price changes over short intervals with an average $R^2$ of about
65\%---$R^2$, the coefficient of determination, being the fraction of the variance of price changes
that the fitted line explains---against about 32\% for the imbalance of \emph{trades} alone; adding
trade imbalance to OFI raises the average only to about 67\%. The slope is close to inversely
proportional to average depth, $\beta \propto 1/\text{depth}$. The fit varies a great deal across
stocks---from above 90\% for some to below 15\% for others---which is itself a warning against
quoting a single number for ``the'' predictability of order flow. The lesson that limit orders and
cancellations carry as much price information as trades is one of the most robust in the field.

\paragraph{Looking deeper than the best level.} Information also sits beyond the best quotes. On
the Australian Stock Exchange, \citet{cao2009} attribute about 22\% of price discovery to the book
beyond the best bid and offer, and find that imbalances deeper in the book predict short-term
returns even after controlling for the autocorrelation of returns, the inside spread and the imbalance
of trades. \emph{Multi-level OFI} computes the
imbalance separately at each of the first $L$ levels, giving a vector rather than a number; for six
Nasdaq stocks, \citet{xu2019mlofi} find that the out-of-sample fit improves with every additional
level included. \citet{contcucuringu2023} combine the level-wise imbalances into a single
\emph{integrated OFI}, the first principal component of the multi-level vector, which captures over
89\% of its variance. On the 100 largest S\&P~500 stocks over 2017--2019 the integrated OFI raises
the average out-of-sample $R^2$ of the contemporaneous price-impact regression from about 65\%
(best level only) to about 84\%.

\paragraph{Order flow in other instruments.} The same study asks whether the order flow of
\emph{other} stocks helps (cross-impact, Section~\ref{sec:cross}). Two answers are needed. For the
\emph{contemporaneous} price change, once integrated OFI is used, adding other stocks' order flow
does not improve the out-of-sample fit. For \emph{forecasting} the next minute's return, lagged
order flow of other stocks does help significantly, although the out-of-sample $R^2$ of all the
forecasting models remains slightly negative and the effect decays within minutes. Explaining the
present and predicting the future are, again, different problems.

\paragraph{Horizon and reversal.} The predictive content of order imbalance depends on the horizon.
In Chinese equities, \citet{zhang2025oi} find that order imbalance predicts returns positively from
5 to 30 minutes ahead and negatively from 60 to 120 minutes ahead, with no reversal for large or
heavily traded stocks. The usual interpretation is that part of the price pressure from imbalance
is temporary and is undone as liquidity is replenished. A signal must be used over the horizon on
which it has the sign one assumes.

\paragraph{Estimating the slope in real time.} Because $\beta$ moves with depth, it is re-estimated
continuously in practice: by a rolling-window regression, by normalising OFI by recent average
depth (a \emph{depth-normalised} OFI), or by treating $\beta_t$ as a hidden state that drifts over
time and tracking it with a Kalman filter~\citep{kalman1960}, the standard recursive estimator for
linear systems observed with noise.

\paragraph{What OFI is used for.} The same quantity serves several decisions. As a directional
signal, it predicts the sign of the next short-horizon move. For a market maker, it sets the
\emph{skew} of quotes---leaning bids and offers in the direction of pressure---and informs inventory
control~\citep{cartea2015}. For execution, it indicates when resting passively is likely to be
rewarded and when the order book is about to move away.

\begin{modelcard}{Order-flow imbalance}
\cardrow{What you calibrate}{The price-impact slope $\beta$ for a chosen window $w$ (and, for
multi-level OFI, one slope per level or the weights of the integrated OFI).}
\cardrow{How}{An ordinary least-squares regression of mid changes on $\OFI^{(w)}_t$. Because $\beta$
falls as depth rises, it is re-estimated on a rolling window, absorbed by dividing OFI by recent
depth, or tracked as a drifting hidden state with a Kalman filter. Integrated OFI takes the first
principal component of the level-wise OFIs.}
\cardrow{What you do with it}{A baseline directional signal; the skew of a market maker's quotes; a
warning, for a passive order, that the book is about to move away from it.}
\cardrow{How to picture it}{Net flow into a tank: lots joining the bid or leaving the ask are inflow,
lots joining the ask or leaving the bid are outflow; the level of the tank leads the price.}
\end{modelcard}

\subsection{Queue imbalance and richer book features}

Queue imbalance compresses the top of book into one number,
\begin{equation}
  \QI_t = \frac{v^b_{1,t} - v^a_{1,t}}{v^b_{1,t} + v^a_{1,t}} \in [-1, 1].
\end{equation}
\inwords{the difference between the best-bid and best-ask queue sizes, divided by their total so
the result does not depend on how busy the market is. It is $+1$ when only buyers are queued, $-1$
when only sellers are, and $0$ when the two queues are equal.}

\paragraph{Why imbalance predicts the next move: the thin side breaks first.}
The mid moves when one of the two best queues is used up. If the book is skewed---many lots on one
side, few on the other---the thin queue needs far less trading and cancelling to disappear than the
thick one, so it is likely to be emptied first. The next price move is therefore more likely to be
\emph{towards the thin side}: a thin ask is likely to be taken out, and the price ticks up; a thin
bid is likely to be hit, and the price ticks down. Think of two piles of sand being removed at
roughly the same rate: the smaller pile is gone first. Stochastic models of the book make this
precise, computing the probability of an up move as a function of the two queue sizes~\citep{cst2010}.

\begin{example}[Reading the imbalance in Figure~\ref{fig:book}]
The best bid has 29 lots and the best ask 14, so
$\QI_t = (29 - 14)/(29 + 14) \approx 0.35$: the book leans towards the buyers. The ask queue is the
thinner one and is likely to break first, so the next change in the mid is more likely to be up
than down. For a trader the skew has two consequences. A new passive buy
order would join the back of the long 29-lot queue and probably wait while the price moves away
from it. A new passive sell order would join the short ask queue and probably be filled---just
before the up move that the imbalance predicted, which is adverse selection
(Section~\ref{sec:adverse}).
\end{example}

\paragraph{A price that leans with the book.} The same intuition gives a better single-number price
than the mid. The \emph{weighted mid} weights each side's price by the size on the \emph{other}
side,
\begin{equation}
  m^{\text{wmid}}_t = \frac{v^b_{1,t}\, p^a_{1,t} + v^a_{1,t}\, p^b_{1,t}}{v^b_{1,t} + v^a_{1,t}},
\end{equation}
\inwords{when the bid queue is much larger than the ask queue, the price is pulled towards the ask,
the side expected to break, and vice versa. With equal queues it equals the mid. In
Figure~\ref{fig:book} it is $5000.00 + 0.25 \times 29/43 \approx 5000.169$, above the mid of
5000.125.}

\begin{modelcard}{Queue imbalance and weighted mid}
\cardrow{What you calibrate}{Nothing for $\QI_t$ or $m^{\text{wmid}}_t$ themselves. To turn $\QI_t$
into a probability, the coefficients of a logistic regression of the next mid-move direction on
$\QI_t$~\citep{gould2016qi}.}
\cardrow{How}{Count, over historical data, which way the mid moved next after each observed value
of the imbalance, and fit the logistic curve by maximum likelihood; report it separately for
large-tick and small-tick instruments, where its value differs.}
\cardrow{What you do with it}{Predict the direction of the next mid change; decide which side a
passive order should rest on and when it is likely to be adversely selected.}
\cardrow{How to picture it}{Two piles of sand removed at about the same rate: the smaller pile is gone
first, and the price moves towards the side whose pile disappeared.}
\end{modelcard}

\subsection{The micro-price}
\label{sec:microprice}

The weighted mid has three defects~\citep{stoikov2018}. It is not a \emph{martingale}---its
expected future value is not its current value, so it is a biased estimate of where the price is
going. It is noisy, especially when the spread widens. And it can move the wrong way: if the last
lot on a thin ask is cancelled, the ask moves up a tick, yet the weighted mid can fall. The
\emph{micro-price} of \citet{stoikov2018} replaces the ad hoc weighting by an estimate of the
expected future mid itself.

\paragraph{Definition.} Let $\eta_1 < \eta_2 < \cdots$ be the successive times after $t$ at which
the mid changes, let $\iota_t = v^b_{1,t}/(v^b_{1,t} + v^a_{1,t}) = (1+\QI_t)/2 \in [0,1]$ be the
imbalance rescaled to $[0,1]$, and let $s_t$ be the spread. The micro-price is the limit of the
expected mid after more and more price changes,
\begin{equation}
  m^{\text{micro}}_t = \lim_{k\to\infty} \E\big[m_{\eta_k} \mid \iota_t, s_t\big].
\end{equation}
\inwords{start from the current imbalance and spread, and ask where the mid will be after the next
price change, after the one after that, and so on. The answers settle down to a value; that value is
the micro-price. By construction it is the best estimate, given the top of the book, of where the
price is heading.}

Two assumptions make it computable: the pair $(\iota_t, s_t)$ evolves as a Markov chain, and the
dynamics do not depend on the price level itself. Then
\begin{equation}
  \begin{gathered}
  m^{\text{micro}}_t = m_t + \sum_{k=1}^{\infty} g^{(k)}(\iota_t, s_t), \\
  g^{(1)}(\iota, s) = \E\big[m_{\eta_1} - m_t \mid \iota, s\big], \qquad
  g^{(k+1)}(\iota, s) = \E\big[g^{(k)}(\iota_{\eta_1}, s_{\eta_1}) \mid \iota, s\big].
  \end{gathered}
\end{equation}
\inwords{the first term is the expected size of the next price change; each further term is the
expected next change, averaged over the states the book can be in after the previous one. The
micro-price is the mid plus the sum of all these expected adjustments.}

\paragraph{Estimation.} In practice the imbalance is split into $n_\iota$ buckets (for example by
quantiles) and the spread into $n_s$ values, so the bucketed state $(\iota, s)$ takes $n_\iota n_s$
values. From data one counts three kinds of transition from each state: to another state with the mid
unchanged, collected in a matrix $\mathbf{U}$; to each possible size of mid change, in a matrix
$\mathbf{J}^{(1)}$; and to another state with the mid changed, in a matrix $\mathbf{J}^{(2)}$. With
$\boldsymbol{\Delta}$ the vector of possible mid changes and $\mathbf{B} = (I -
\mathbf{U})^{-1}\mathbf{J}^{(2)}$,
\[
  g^{(1)} = (I - \mathbf{U})^{-1} \mathbf{J}^{(1)} \boldsymbol{\Delta}, \qquad
  g^{(k+1)} = \mathbf{B}\, g^{(k)}.
\]
The factor $(I-\mathbf{U})^{-1}$ sums over all the steps in which the book changes without
the mid moving---the same geometric-series device used for the Hawkes model in Section~\ref{sec:hawkes}. The data
are symmetrised before estimation (each observation is added again with bid and ask swapped and the
price change negated), which guarantees that the series converges. Stoikov reports that a handful
of terms suffices in practice, and that the micro-price predicts short-term price changes better
than either the mid or the weighted mid, for both a large-tick and a small-tick stock.

\begin{modelcard}{Micro-price}
\cardrow{What you calibrate}{The transition matrices $\mathbf{U}$, $\mathbf{J}^{(1)}$,
$\mathbf{J}^{(2)}$ between imbalance--spread buckets, after choosing the numbers of buckets $n_\iota$
and $n_s$.}
\cardrow{How}{Bucket the imbalance and spread, symmetrise the data, and count transitions; the
adjustments $g^{(k)}$ then follow by matrix algebra, and a handful of terms suffices.}
\cardrow{What you do with it}{A fair-value estimate to quote around in place of the mid, and a
better reference than the mid or weighted mid for predicting short-term price changes.}
\cardrow{How to picture it}{A weather forecast that starts from today's state and averages over
every path the next few changes can take: the micro-price is where those paths land on average.}
\end{modelcard}

\subsection{The shape of the book}
\label{sec:shape}

Beyond the best level, several features summarise how liquidity is distributed:
\begin{description}[leftmargin=1.5em,style=nextline,itemsep=2pt]
\item[Multi-level imbalance] the queue imbalance computed with sizes summed over the best $\ell$
  levels on each side, for several $\ell$.
\item[Book slope] how fast cumulative depth grows with distance from the best price, usually the
  fitted slope of cumulative size against price distance, separately for each side. A steep slope
  means a deep book that absorbs large orders.
\item[Curvature] the change of that slope with distance---whether depth is concentrated near the
  touch or further out. Its empirical relevance is tied to the average shape of the book, which in
  equities is humped: depth is usually largest a few ticks away from the best price, not at
  it~\citep{bouchaud2002}.
\item[Depth-weighted and volume-weighted prices] averages of the prices at the first $\ell$ levels,
  weighted by their sizes; the weighted mid is the one-level case.
\item[Resilience] how quickly the book refills after a large trade. \citet{large2007} measures it with
  a ten-variate Hawkes model on an LSE stock and finds that when the book does refill, it does so
  the impulse responses have half-lives under 20 seconds, but a resilient response follows a large
  trade in under 40\% of cases.
\item[Depth-normalised quantities] order-flow measures divided by recent average depth, motivated by
  the inverse relation between price impact and depth.
\end{description}

\begin{table}[ht]
\centering\footnotesize
\caption{Classical order-book signals at a glance. Horizons are indicative of how the signals are
used in the cited literature, not guarantees.}
\label{tab:signals}
\setlength{\tabcolsep}{3pt}
\begin{tabularx}{\linewidth}{l*{4}{>{\raggedright\arraybackslash}X}}
\toprule
Signal & What it measures & Typical horizon & How it is fitted & Known failure modes \\
\midrule
OFI, best level~\citep{cks2014} & net buying pressure at the touch & seconds to minutes (contemporaneous fit) & linear regression; slope re-estimated with depth & weak in thin, noisy names; distorted by spoofing \\
Multi-level / integrated OFI~\citep{xu2019mlofi,contcucuringu2023} & pressure across the first levels & same & regression; principal components & more inputs to monitor; depth beyond the touch is less committed \\
Queue imbalance~\citep{gould2016qi} & skew of the two best queues & next mid change & logistic regression & weak for small-tick names; misled by hidden size \\
Weighted mid & price leaning toward the thin side & next mid change & none & noisy; not a martingale \\
Micro-price~\citep{stoikov2018} & expected future mid given imbalance and spread & a few mid changes & Markov chain on imbalance and spread buckets & assumes Markov dynamics; needs symmetrised data \\
Book slope, curvature & how liquidity is spread across levels & longer & descriptive & sensitive to transient deep orders \\
Resilience~\citep{large2007} & speed of refill after trades & seconds & Hawkes impulse response & varies strongly across regimes \\
\bottomrule
\end{tabularx}
\end{table}

\paragraph{Limits of the intuition.} Displayed size is not committed size. A thin queue may be
refilled immediately by new orders; hidden or iceberg orders may make the thin side thicker than it
looks; and orders placed only to be cancelled (spoofing, Section~\ref{sec:failures}) can create a
skew that is not real. Imbalance is therefore a strong but noisy signal.

\paragraph{Evidence.} \citet{gould2016qi} fit a logistic regression of the direction of the next
mid move on queue imbalance for ten Nasdaq stocks. The relationship is strongly significant for
every stock, but its value depends on the tick size: out of sample, queue imbalance improves binary
classification over a null model by about 50--60\% for large-tick stocks and only about 10--30\% for
small-tick stocks, and probabilistic classification by about 20--30\% against about 2--6\%. Queue
imbalance is most informative exactly where queues are long and the price rarely moves---the
large-tick regime that ES and ZN futures also belong to.

\begin{example}[Same imbalance, different history]
\label{ex:qi}
Bid 300, ask 100 gives $\QI_t = (300-100)/(300+100) = 0.5$. In book~A the ask fell from 400 to 100
because a 300-lot market buy lifted it. In book~B it fell because 300 lots were cancelled. Both
books have the same snapshot and the same $\QI_t$. A trader who saw the messages knows that A
contains an aggressive buyer and B does not.
\end{example}

\begin{modelcard}{Book shape}
\cardrow{What you calibrate}{The number of levels $\ell$ to aggregate over; the slope and curvature
of cumulative depth against distance; for resilience, the impulse response of a Hawkes model.}
\cardrow{How}{Slope and curvature are least-squares fits of cumulative size against price distance,
per side and per snapshot. Resilience is read off a fitted multivariate Hawkes model
(Section~\ref{sec:hawkes}), as in~\citet{large2007}.}
\cardrow{What you do with it}{Estimate how much a large order will move the price, normalise other
signals by depth, and detect a book that is thinning out before the touch does.}
\cardrow{How to picture it}{The profile of a hillside on each side of a valley: a steep hillside
holds a lot of liquidity close to the price, a shallow one makes large orders walk far.}
\end{modelcard}

The central empirical question is not whether richer representations improve in-sample fit, but
whether the improvement survives out-of-sample evaluation and execution costs.

% ================================================================
\section{Events rather than snapshots}
\label{sec:events}
% ================================================================

The same reduction in displayed volume has different information content depending on its cause.
A snapshot representation can make two states identical even when their histories differ
(Example~\ref{ex:qi}). Market-by-order (L3) data carry the cause: every add, modify, cancel and
execution, with an order identifier. On CME's MDP~3.0 feed, for instance, book updates for
individual orders and trade summaries arrive as separate message templates, so the two books in
Example~\ref{ex:qi} produce visibly different message sequences.

Treating the three event types jointly is not optional. Price changes are often attributed to market
orders, but in Nasdaq data roughly half of all price-changing events are limit orders placed
\emph{inside} the spread, which improve the best price rather than consume it~\citep{eisler2012}.
A model that sees only trades misses half of what moves the price.

The case for message-level representation has recently been made directly in deep learning.
LOBERT~\citep{linna2025lobert} adapts the BERT encoder to LOB data by treating each complete
multi-dimensional message as a single token, while keeping continuous representations of price,
volume and time. It is designed as a general-purpose, fine-tunable foundation model, and reports
leading performance on mid-price movement and next-message prediction with a shorter required
context than earlier methods; it is trained and evaluated on Nasdaq ITCH message data for four
large US stocks in 2015, with millisecond timestamps. A message-level encoder sees the difference between a cancellation
and an execution in Example~\ref{ex:qi}; a snapshot encoder, by construction, does not.

% ================================================================
\section{Queue depletion and first-passage models}
\label{sec:queue-models}
% ================================================================

\paragraph{The idea: a queue is a bucket with a leak.} Think of the best-ask queue as a bucket of
lots. New sell limit orders pour lots in; buy market orders that trade against it, and sellers who
cancel, drain lots out. Each event adds or removes a random amount at a random time, so the level of
the bucket wanders up and down like a random walk. The question this section asks is: \emph{how long
until the bucket is empty?} That moment matters for three reasons.
\begin{enumerate}[leftmargin=1.8em,itemsep=1pt]
\item When the best-ask queue empties, the best ask moves up a tick, and so does the mid. Price
  changes \emph{are} queue depletions.
\item The direction of the next price move is a race between the two best queues: whichever empties
  first decides it. This is the formal version of ``the thin side breaks first''
  (Section~\ref{sec:ofi}).
\item For a resting order, the wait until it is filled is the time until the volume \emph{ahead of
  it} is used up---the same question, asked about part of a queue (Section~\ref{sec:queue-position}).
\end{enumerate}

\paragraph{The definition.} The depletion time of the best ask is the first-passage time
\begin{equation}
  \tau^a = \inf\{u > 0 : v^{a}_{1,t+u} \le 0\},
\end{equation}
where $u$ runs over time after $t$; $\tau^b$ is defined in the same way for the bid.
\inwords{$\tau^a$ is how long from now until the best-ask queue is empty for the first time---the
time at which a loop that polls the queue size would first see zero. ``First passage'' is the
general name for the first time a wandering quantity reaches a level.}
The race between the two queues is then
\begin{equation}
  \P(\text{next mid move is up} \mid X_t) = \P(\tau^a < \tau^b \mid X_t),
\end{equation}
\inwords{the price goes up next exactly when the ask queue empties before the bid queue does.}

\paragraph{A first estimate.} If lots drain out at an average rate of $r_{\text{out}}$ lots per
second and pour in at $r_{\text{in}}$, the bucket loses $r_{\text{out}} - r_{\text{in}}$ lots per
second on average, and a best-ask queue of $v^a_{1,t}$ lots takes on average
\[
  \E[\tau^a] \approx \frac{v}{r_{\text{out}} - r_{\text{in}}}
  \qquad (r_{\text{out}} > r_{\text{in}})
\]
seconds to empty: size divided by net drain rate, as for any bucket. Two warnings come with this
formula. First, it is an average; the actual time is random and often far from it. Second, if
inflow keeps up with outflow ($r_{\text{in}} \ge r_{\text{out}}$), the bucket has no net leak and the
expected time can be very long or infinite.

\paragraph{Models.} Stochastic models of the book~\citep{cst2010} treat arrivals, cancellations
and executions as random event streams (point processes) and derive the distribution of
first-passage times and the probability of a mid move in either direction, including the race above.
Queue-reactive models~\citep{huang2015} go one step further and let the rates themselves depend on
the current queue sizes---for example, orders may join or cancel at different rates when a queue is
short than when it is long. That reproduces many observed regularities of large-tick assets.

\subsection{The Cont--Stoikov--Talreja model}
\label{sec:cst}

The model of \citet{cst2010} is the foundation on which most later queue-based models of the
order book are built, and it repays careful study. It was one of the first models of the whole book to
satisfy three requirements at once: its parameters can be estimated directly from order-book
data with simple counting; it reproduces important empirical features of the book without being tuned
to do so; and it is simple enough that the probabilities a trader actually needs---will the price go
up next, will my order be filled before the price moves---can be \emph{computed}, quickly and exactly,
rather than estimated by simulation. Its central idea is that a limit order book is a queueing
system: limit orders wait in line to be served by market orders, or leave the line by cancelling.

\subsubsection{The model}

\paragraph{State.} Prices live on a grid of ticks. The state of the book is the number of orders
resting at every price on the grid, counted in units of the average limit-order size. Because most
activity happens near the current price, the model measures position relative to the quotes: let
$\mathcal{Q}^{\text{bid}}_{\bar\ell}(t)$ be the number of buy orders $\bar\ell$ ticks below the best
ask, and $\mathcal{Q}^{\text{ask}}_{\bar\ell}(t)$ the number of sell orders $\bar\ell$ ticks above the
best bid. Note the convention: distance is measured from the \emph{opposite} best quote. With a
one-tick spread, the best bid is at distance $\bar\ell = 1$ from the best ask.

\paragraph{Events and their rates.} Six kinds of event change the book, and each is modelled as an
independent Poisson stream (Section~\ref{sec:hawkes} explains what that assumption means and where
it fails):
\begin{itemize}[leftmargin=1.5em,itemsep=2pt]
\item a \textbf{limit buy} order arrives $\bar\ell$ ticks below the best ask, and a \textbf{limit sell}
  order $\bar\ell$ ticks above the best bid, each at rate $r_L(\bar\ell)$;
\item a \textbf{market buy} order removes one unit from the best ask, and a \textbf{market sell} order
  one unit from the best bid, each at rate $r_M$;
\item each resting order $\bar\ell$ ticks from the opposite quote is \textbf{cancelled} independently at
  rate $r_c(\bar\ell)$, so a level holding $\mathcal{Q}$ orders loses orders to cancellation at total
  rate $r_c(\bar\ell)\,\mathcal{Q}$.
\end{itemize}
Two modelling choices carry the economics. First, limit orders arrive more often near the current
price; following earlier empirical work~\citep{bouchaud2002}, the arrival rate is a power law in
distance,
\begin{equation}
  r_L(\bar\ell) = \frac{r_0}{\bar\ell^{\,\Xi}} ,
\end{equation}
with a scale $r_0$ and exponent $\Xi$ fitted to data.
\inwords{the farther a price is from where trading happens, the fewer orders are placed there, and the
fall-off follows a power law: doubling the distance divides the rate by $2^{\Xi}$.}
Second, cancellations are proportional to queue size: each order has the same chance per second of
being cancelled, so a long queue sheds orders faster than a short one. This single assumption is what
keeps queues from growing without bound.

Under these assumptions the book is a continuous-time Markov chain: its future depends only on the
current queue sizes, not on how they were reached. The authors show that a book that starts with the
bid below the ask always keeps that property, and that if every cancellation rate is positive the
chain is \emph{ergodic}---it has a unique long-run distribution, so averages over a long simulated
history converge to well-defined values, such as the average shape of the book.

\subsubsection{Estimating the model}

The estimation is deliberately elementary, which is part of the model's appeal. With the unit of size
set to the average limit order:
\begin{itemize}[leftmargin=1.5em,itemsep=2pt]
\item the limit-order rate at each distance is the number of times the queue at that distance
  \emph{increased}, divided by the trading time, $\widehat{r}_L(\bar\ell) = N_L(\bar\ell)/T_{\text{obs}}$,
  and a power law is fitted to the first five distances by least squares to extrapolate further out;
\item the market-order rate is the number of market orders that changed the best quotes, scaled by
  the ratio of average market-order size to average limit-order size, divided by the trading time;
\item the per-order cancellation rate at each distance is the number of decreases not caused by
  trades, scaled by size, divided by the trading time \emph{and by the average queue size at that
  distance}---because the model makes total cancellations proportional to queue size.
\end{itemize}
Here $N_L(\bar\ell)$ is the count of queue increases at distance $\bar\ell$ and $T_{\text{obs}}$ the
total observed trading time. The authors use five levels of quotes and trades for Tokyo Stock Exchange
stocks over 125 trading days (August to December 2006) and report estimates for Sky Perfect
Communications (Table~\ref{tab:cst-est}).

\begin{table}[ht]
\centering\small
\caption{Estimated Cont--Stoikov--Talreja parameters for Sky Perfect Communications, Tokyo Stock
Exchange, 2006~\citep{cst2010}. Rates are per minute, in units of the average limit-order size.}
\label{tab:cst-est}
\begin{tabular}{lccccc}
\toprule
Distance from opposite quote $\bar\ell$ & 1 & 2 & 3 & 4 & 5 \\
\midrule
Limit-order rate $\widehat{r}_L(\bar\ell)$ & 1.85 & 1.51 & 1.09 & 0.88 & 0.77 \\
Per-order cancellation rate $\widehat{r}_c(\bar\ell)$ & 0.71 & 0.81 & 0.68 & 0.56 & 0.47 \\
\midrule
\multicolumn{6}{l}{Market-order rate $\widehat{r}_M = 0.94$; power-law fit $\widehat{r}_0 = 1.92$, $\widehat{\Xi} = 0.52$} \\
\bottomrule
\end{tabular}
\end{table}

\subsubsection{The key simplification: two independent birth--death queues}

The full book is a large Markov chain, but the questions a trader asks concern only the next price
change, and until the price changes the quotes do not move. During that time the best-bid queue and
the best-ask queue each evolve as a \emph{birth--death process}: a counter that goes up by one at a
constant ``birth'' rate and down by one at a ``death'' rate that depends on its current value. If the
spread is $s_t/\delta$ ticks, the best queues sit at distance $\bar\ell = s_t/\delta$ from the opposite
quote, and a queue holding $x$ orders evolves as
\begin{center}
\begin{tikzpicture}[
  st/.style={draw, circle, minimum size=8mm, font=\small},
  arr/.style={-{Stealth[length=2mm]}, thick}]
\node[st, fill=gray!15] (s0) at (0,0) {0};
\node[st] (s1) at (2.4,0) {1};
\node[st] (s2) at (4.8,0) {2};
\node[st] (s3) at (7.2,0) {3};
\node at (9.0,0) {$\cdots$};
\draw[arr, bend left=30] (s2) to node[above, font=\scriptsize]{$r_L$} (s3);
\draw[arr, bend left=30] (s1) to node[above, font=\scriptsize]{$r_L$} (s2);
\draw[arr, bend left=30] (s3) to node[below, font=\scriptsize]{$r_M + 3r_c$} (s2);
\draw[arr, bend left=30] (s2) to node[below, font=\scriptsize]{$r_M + 2r_c$} (s1);
\draw[arr, bend left=30] (s1) to node[below, font=\scriptsize]{$r_M + r_c$} (s0);
\end{tikzpicture}
\end{center}
with birth rate $r_L(s_t/\delta)$ and death rate $r_M + x\,r_c(s_t/\delta)$; reaching state $0$ (shaded) means
the queue has been used up and the price moves, so the count stops there. The authors prove that, up to the first price change,
the two best queues can be replaced by \emph{independent} birth--death processes with these rates. Every
question about the next price move becomes a question about which of two independent counters hits
zero first---the race of Section~\ref{sec:queue-models}, now with exact rates.

\paragraph{Why Laplace transforms.} The time for a birth--death process starting at $x$ to reach zero is
a sum of $x$ independent pieces: the time to go from $x$ to $x-1$, then from $x-1$ to $x-2$, and so on.
Distributions of sums are awkward, but the \emph{Laplace transform} of a waiting time $\tau$,
$\E[\mathrm{e}^{-\varpi \tau}]$ as a function of a variable $\varpi$, turns sums of independent waiting
times into products of their transforms. For birth--death processes each one-step transform has a known
closed form as a continued fraction~\citep{cst2010}, so the transform of the whole first-passage time is a
product of known pieces. Probabilities such as ``queue A empties before queue B'' are then obtained by
numerically inverting a transform built from these products---a computation that takes far less time than
simulating thousands of paths of the book, and gives full distributions rather than only averages.
\inwords{instead of simulating the book many times and counting how often the price went up, the model
writes the answer as a formula in transform space and evaluates it directly. The transform is a
book-keeping device that makes adding independent waiting times as easy as multiplying numbers.}

\subsubsection{Three probabilities a trader needs}

The paper computes three conditional probabilities, each given the current number of orders at the
best bid, $v^b_{1,t}$, and at the best ask, $v^a_{1,t}$, and the spread.

\paragraph{1. Will the next price move be up?} With a one-tick spread, the mid rises next if the ask
queue empties before the bid queue. With a wider spread, a new order can also arrive inside the spread
and move the mid first; the formula adds these arrivals as competing exponential clocks.
Table~\ref{tab:cst-up} compares the model's answer with the empirical frequency in the data. The model
captures the main pattern well---the probability rises with the bid queue and falls with the ask
queue---though it understates how strongly a large ask queue pushes the price down, and some cells
differ noticeably (for example 0.548 against an empirical 0.714 with four orders at the ask and five
at the bid).

\begin{table}[ht]
\centering\footnotesize
\caption{Probability that the next mid move is up, one-tick spread, Sky Perfect Communications:
empirical frequency / Cont--Stoikov--Talreja model (Laplace-transform method)~\citep{cst2010}. Queue sizes
in units of the average limit order.}
\label{tab:cst-up}
\begin{tabular}{lccccc}
\toprule
& \multicolumn{5}{c}{Orders at the best ask, $v^a_{1,t}$} \\
\cmidrule(lr){2-6}
Orders at the best bid, $v^b_{1,t}$ & 1 & 2 & 3 & 4 & 5 \\
\midrule
1 & 0.512 / 0.500 & 0.304 / 0.336 & 0.263 / 0.259 & 0.242 / 0.216 & 0.226 / 0.188 \\
2 & 0.691 / 0.664 & 0.502 / 0.500 & 0.444 / 0.407 & 0.376 / 0.348 & 0.359 / 0.307 \\
3 & 0.757 / 0.741 & 0.601 / 0.593 & 0.533 / 0.500 & 0.472 / 0.437 & 0.409 / 0.391 \\
4 & 0.806 / 0.784 & 0.672 / 0.652 & 0.580 / 0.563 & 0.529 / 0.500 & 0.484 / 0.452 \\
5 & 0.822 / 0.812 & 0.731 / 0.693 & 0.640 / 0.609 & 0.714 / 0.548 & 0.606 / 0.500 \\
\bottomrule
\end{tabular}
\end{table}

\paragraph{2. Will my order at the bid be filled before the price moves?} Suppose an order is placed at
the back of the best bid, so that the bid queue holds $v^b_{1,t}$ orders including our own, and that it
is never cancelled. It is filled when all orders ahead of it
have been removed---by market sells, or by cancellation of the orders ahead---which is a pure-death
process: with $x$ orders in line including our own, the line shortens at rate $r_M + (x-1)\,r_c$, since
our own order is not cancelled. The order is filled before the price moves if this line empties before
the ask queue does. Table~\ref{tab:cst-fill} gives the model's answers. The pattern is the execution
problem of Part~III in numbers: an order alone at the bid, facing five orders at the ask, is filled before
the price moves about 78\% of the time; the last of five orders at the bid, facing one at the ask, about
12\%.
The authors note that their data do not identify individual orders, so these probabilities could not be
checked against data---exactly the gap that order-level data and queue-exact replay fill
(Part~IV).

\begin{table}[ht]
\centering\footnotesize
\caption{Probability that an order placed at the best bid is executed before the mid moves, one-tick
spread, Cont--Stoikov--Talreja model (Laplace-transform method)~\citep{cst2010}.}
\label{tab:cst-fill}
\begin{tabular}{lccccc}
\toprule
& \multicolumn{5}{c}{Orders at the best ask, $v^a_{1,t}$} \\
\cmidrule(lr){2-6}
Orders at the best bid, $v^b_{1,t}$ & 1 & 2 & 3 & 4 & 5 \\
\midrule
1 & 0.497 & 0.641 & 0.709 & 0.749 & 0.776 \\
2 & 0.302 & 0.449 & 0.535 & 0.591 & 0.631 \\
3 & 0.206 & 0.336 & 0.422 & 0.483 & 0.528 \\
4 & 0.152 & 0.263 & 0.344 & 0.404 & 0.452 \\
5 & 0.118 & 0.213 & 0.287 & 0.346 & 0.393 \\
\bottomrule
\end{tabular}
\end{table}

\paragraph{3. Will I ``make the spread''?} A market maker places one order at the best bid and one at
the best ask and earns the spread if both are filled before the price moves. The model computes this
probability too. Its most interesting feature is that it is \emph{not} monotone: for a fixed number of
orders at the ask, the probability first rises and then falls as the bid queue grows. With one order at
the ask, it is about 0.27 with one order at the bid, peaks near 0.31 with two or three, and falls back
to about 0.29 with five. A longer bid queue protects the bid from being exhausted, which helps, but also
makes the bid order wait longer, which hurts; the balance gives an interior optimum.

\subsubsection{What the model reproduces}

Without any tuning for these purposes, the fitted model reproduces several features of the real book:
\begin{itemize}[leftmargin=1.5em,itemsep=2pt]
\item the \textbf{hump-shaped average book}: simulated average depth peaks two ticks from the opposite
  quote, as observed empirically~\citep{bouchaud2002};
\item the \textbf{order of magnitude of volatility}: over the first day of the sample the realised
  volatility was 0.0219 after 370 trades, against $0.0228 \pm 0.0003$ in simulations of the same number
  of trades---computed from order-flow rates alone;
\item the \textbf{one-step behaviour of queues}: the probability that a queue grows at its next change,
  as a function of its size and distance, is close to the empirical frequencies at all five distances.
\end{itemize}

\subsubsection{A trading test, and what it shows}

The authors build a simple strategy on the probability of an up move. When the spread is one tick, there
are at least three orders at the best bid, one at the best ask and at most one at the second ask, buy
with a market order. If one order remained at the second ask, the book now has a two-tick spread, three
or more orders at the bid and one at the ask, and the model gives a 0.62 probability that the next move
is up. Exit with a market sell when the bid rises above the entry price, or when the bid queue falls to
one order (a one-tick loss). In about 30 days of \emph{simulated} trading from the fitted model, the
strategy made 2,376 round trips with an average profit of 0.068 ticks per trade, although fewer than
half of the trades were winners---and without transaction costs.

The test illustrates both the promise and the limits discussed throughout this survey. The decision is
derived from a conditional probability of the book, which is exactly what an execution policy needs.
But it is evaluated inside the model that generated the probability, not on recorded data, and without
costs; a profit of 0.068 ticks per trade must also survive the half spread paid on entry and exit, the
fees, and the latency between seeing the book and the market order arriving. These are the steps of the
chain in Equation~(\ref{eq:chain}) that the model does not cover.

\subsubsection{Limitations and descendants}

The authors are explicit about what the model leaves out. Order arrivals are independent Poisson
streams, so there is no clustering of events and none of the long memory of order flow
(Section~\ref{sec:facts}); orders all have the same size; rates do not change through the day or with
market conditions; and there is no strategic behaviour. Each limitation defines a line of later work:
\begin{itemize}[leftmargin=1.5em,itemsep=2pt]
\item \citet{contlarrard2013,contlarrard2012} reduce the book to the two best queues and obtain closed
  forms and a diffusion limit (Section~\ref{sec:cdl});
\item queue-reactive models let all rates depend on the current queue sizes~\citep{huang2015}, and their
  neural extensions learn that dependence~\citep{deepqr};
\item Hawkes models replace the Poisson streams by self-exciting ones (Section~\ref{sec:hawkes});
\item the micro-price~\citep{stoikov2018} keeps the Markov-chain logic but estimates the transitions
  directly from data rather than from a parametric model;
\item queue-position valuation~\citep{moallemi2017} and fill-probability models~\citep{arroyo2024,kanformer}
  extend the second probability above to individual orders, with adverse selection.
\end{itemize}
Finally, because the Laplace-transform method yields full distributions rather than single
probabilities, it can answer latency questions directly---for example, the probability that the ask
queue empties more than a given delay after the bid queue, which is what matters when an order takes
time to reach the exchange~\citep{cst2010}.

\begin{modelcard}{Cont--Stoikov--Talreja}
\cardrow{What you calibrate}{The limit-order rates $r_L(\bar\ell)$, summarised by $r_0$ and $\Xi$; the
market-order rate $r_M$; the per-order cancellation rates $r_c(\bar\ell)$; the unit of size (the
average limit order).}
\cardrow{How}{Counting: increases, decreases not caused by trades, and market orders at each distance,
divided by trading time (and, for cancellations, by average queue size); a least-squares power-law fit
for $r_L$.}
\cardrow{What you do with it}{Compute, from the current two best queues, the probability that the next
move is up, that a resting order fills before the price moves, and that a market maker earns the
spread---fast and exactly, via Laplace transforms.}
\cardrow{How to picture it}{Two counters racing to zero, each ticking up when an order joins and down
when one is traded or cancelled; the first to hit zero decides the next price move.}
\end{modelcard}

\subsection{The Cont--de Larrard model}
\label{sec:cdl}

\citet{contlarrard2013} give the cleanest closed-form treatment of the race between the two best
queues, and it is worth stating because every term has a direct trading meaning. The model assumes:
\begin{itemize}[leftmargin=1.5em,itemsep=1pt]
\item the spread is always one tick (true over 98\% of the time for the large-tick US stocks they
  study);
\item at each best queue, limit orders arrive at rate $r_L$, market orders at rate $r_M$ and
  cancellations at rate $r_C$, all as independent Poisson streams of unit-size orders, so each queue
  is a random walk that moves up at rate $r_L$ and down at rate $r_M + r_C$
  (in the notation of the bucket picture, $r_{\text{in}} = r_L$ and $r_{\text{out}} = r_M + r_C$);
\item when a queue empties, the price moves one tick and both new best queues are drawn afresh from a
  distribution estimated from the data.
\end{itemize}
Empirically the rates nearly balance: $r_L$ is slightly below $r_M + r_C$ (for example a ratio of
about $0.95$ for Citigroup in their data). Three results follow.

\paragraph{Time to the next price change.} The waiting time has a closed-form distribution in terms
of Bessel functions. Its tail is the important part: when inflow is slightly smaller than outflow,
the probability that no price change has occurred after $u$ seconds decays like $1/u^2$; when they
balance exactly, it decays like $1/u$, so the \emph{expected} waiting time is infinite. Long
quiet periods punctuated by bursts are thus a direct consequence of (nearly) balanced queues, before
any self-excitation is added.

\paragraph{Probability that the next move is up.} In the balanced case the probability depends only
on the two queue sizes, not on the rates---it is the probability that a symmetric two-dimensional
random walk started at the queue sizes hits the ask axis first. In the companion heavy-traffic
analysis~\citep{contlarrard2012}, with uncorrelated queue changes of equal size and intensity on the
two sides, it takes the simple form
\begin{equation}
  \P(\text{next mid move is up} \mid v^b_{1,t}, v^a_{1,t}) =
  \frac{2}{\pi}\arctan\!\left(\frac{v^b_{1,t}}{v^a_{1,t}}\right).
\end{equation}
\inwords{the probability rises from $0$ to $1$ as the bid queue grows relative to the ask queue,
and equals $1/2$ when they are equal. It is the ``thin side breaks first'' rule of
Section~\ref{sec:ofi}, made exact for this model. In Figure~\ref{fig:book}, with 29 lots bid and 14
offered, it gives $(2/\pi)\arctan(29/14) \approx 0.71$.} (The general formula in the paper allows
for correlated queue changes; its printed zero-correlation special case has the two queues
interchanged, and the form above is the one implied by the paper's general expression.)

\paragraph{From queues to volatility.} Over long horizons the price behaves like a random walk whose
volatility, in the balanced case, is proportional to $\delta\sqrt{r_L / \Theta}$, where $\Theta$ is the average
product of the bid and ask queue sizes just after a price change. Volatility rises with order
activity and falls with the depth that reappears after each move---a microstructural explanation of
why thin, busy markets are volatile. Across Dow Jones stocks the authors find ten-minute volatility
roughly proportional to this quantity.

\begin{example}[Why the rates must depend on the state]
\label{ex:fpt}
A best-ask queue holds 200 lots. Averaged over the day, net depletion (executions plus
cancellations minus additions) runs at 40 lots per second, so the first estimate gives
$200/40 = 5$ seconds. Now condition the rate on OFI. When buying pressure is in its top decile, the
net drain is 90 lots per second and $\E[\tau^a] \approx 200/90 \approx 2.2$ s; in the bottom decile
it is 15 lots per second and $\E[\tau^a] \approx 13$ s. The same queue has expected lifetimes six
times apart, depending on what the rest of the market is doing. A model that uses one daily average
rate would be wrong in both directions.
\end{example}

\paragraph{Depletion and adverse selection.} Queue depletion is also where adverse selection
comes from. Consider a resting buy order at the back of the best-bid queue. It is filled only when
everything ahead of it has been traded or cancelled. Because it is last in line, in practice it is
filled mostly by a burst of selling large enough to empty the whole level. And an emptied bid level is, by point~1 above, a downward price move. The
order is therefore filled at the very moment the price is moving against it: the buyer acquires
the lots just as they become cheaper. The fill and the bad price move are not two unlucky events
that happen to coincide; they are the same event.

An order near the \emph{front} of the queue is in a different position. It can be filled by a small
sell order that takes a few lots and leaves the rest of the queue, and the price, where they were.
Such fills carry little information, so front-of-queue orders are less adversely selected. This is
one source of the front--back value difference in \citet{moallemi2017}, whose model splits the value
of an order into a static component---the spread earned minus adverse selection, which increases with
queue position---and a dynamic component, the option value of holding the queue position.

\begin{example}[Filled by the move]
In Figure~\ref{fig:book} we join the best bid at 5000.00 with one lot, behind 29 lots. Suppose a
30-lot sell market order arrives. It takes all 29 lots ahead of us and our lot, and the 5000.00
level is gone: the best bid is now 4999.75 and the mid has fallen from 5000.125 to 5000.00 or below.
We bought at 5000.00 and the market immediately values the contract lower. Had we been first in the
queue, a 2-lot sell would have filled us and left 28 lots and the price untouched.
\end{example}

\begin{modelcard}{Cont--de Larrard}
\cardrow{What you calibrate}{Three rates per best queue, $r_L$, $r_M$, $r_C$, and the distribution of
the new best-queue sizes after a price change (which gives $\Theta$).}
\cardrow{How}{Counting events at the best quotes over a sample, and the empirical distribution of queue
sizes just after each price change.}
\cardrow{What you do with it}{The probability that the next move is up, the distribution of time to
the next price change, and a volatility estimate $\propto \delta\sqrt{r_L/\Theta}$ from order-flow rates
alone.}
\cardrow{How to picture it}{A drunkard's walk in two dimensions, starting at the point (bid size, ask
size): whichever axis the walk hits first decides the price move.}
\end{modelcard}

\subsection{The queue-reactive model}
\label{sec:qr}

\paragraph{The idea.} Cont--Stoikov--Talreja make cancellations proportional to queue size and keep
the other rates fixed. \citet{huang2015} ask the data instead: \emph{how do traders actually behave when
a queue is short, and when it is long?} Their answer is the queue-reactive model, in which the rate of
every kind of event at a queue is an arbitrary function of the current size of that queue (and, in
richer versions, of the neighbouring queues), estimated without assuming a functional form.

\paragraph{State and rates.} The book is described by the queues at the first few price levels on each
side of a reference price $p_{\text{ref}}$ (the mid when the spread is an odd number of ticks; three
levels in the paper, because levels four and five behave like level three). Queue sizes are counted in units of
the average size of events at that level, so a queue of size $x$ holds about $x$ typical orders. In the
basic version each queue is a birth--death process like those of Section~\ref{sec:cst}, but with
size-dependent rates: limit orders arrive at rate $r_L(x)$, and cancellations and market orders remove
units at rates $r_C(x)$ and $r_M(x)$, where $x$ is the current queue size.
\inwords{the model writes down a separate arrival, cancellation and trading rate for each possible queue
size, rather than one rate for all sizes. The data decide how the rates change with size.}

\paragraph{Calibration by counting.} For every event at a queue, record the time since the previous
event at that queue, the event type, and the queue size just before the event. The estimator of the
paper is
\[
  \widehat{r}_L(x) = \frac{\#\{\text{limit orders that arrived while the queue held } x\}}
  {\text{total time the queue spent at size } x},
\]
and likewise for cancellations and market orders. (The paper writes it as the inverse of the mean
waiting time at size $x$, times the fraction of events at size $x$ that were limit orders; the two
forms are algebraically equal.) Bid and ask are pooled by symmetry, and recording restarts whenever
$p_{\text{ref}}$ moves.
\inwords{to learn how fast orders arrive when a queue holds $x$ units, count the arrivals that happened
while it held $x$ units and divide by how long it held $x$ units in total.}

\paragraph{What the data show.} On France Telecom and Alcatel-Lucent (Euronext Paris, 2010--2012, five
levels), the rates are far from those of the simpler models. At the best level, the limit-order rate
is roughly constant in queue size (lower when the queue is empty); the cancellation rate increases
concavely up to about 25 units and then flattens, rather than growing linearly as
Cont--Stoikov--Talreja assume; and the market-order rate \emph{decreases} exponentially with the
available volume---traders hit small queues. The long-run distribution of queue sizes implied by the
fitted rates matches the empirical one closely, and better than a Poisson model with linear
cancellation. Richer versions let the best-level rates depend on the opposite queue: insertion falls
and market orders rise as the opposite queue grows.

\paragraph{Prices.} The full queue-reactive model adds a price. When a best queue empties, the
reference price moves one tick with probability $\pi_{\text{move}}$; and with probability $\pi_{\text{reinit}}$
the whole book is redrawn from its long-run distribution around the new price, which the authors
interpret as the share of price changes driven by outside information. ($\pi_{\text{reinit}} = 1$
recovers the Cont--de Larrard assumption of Section~\ref{sec:cdl}.) The two probabilities are calibrated
to the ten-minute volatility and to the ratio of price continuations to reversals; for France Telecom
$\pi_{\text{move}} = 0.7$ and $\pi_{\text{reinit}} = 0.85$. Order-book dynamics alone, with no outside
information, generate at most about 5 basis points of volatility against 14 observed, and far too much
mean reversion; the authors use this gap to motivate the re-initialisation step, which stands for
information arriving from outside the book.

\paragraph{Uses and descendants.} The model is a simulator of large-tick books that reacts to the
orders a trader adds: the paper uses it to compute execution probabilities and to compare placement
tactics (``fire and forget'' against pegging) for a large order. Its neural extension learns the
dependence of the rates on the book with networks~\citep{deepqr}, and queue-reactive Hawkes models add
self-excitation to the size-dependent rates (Section~\ref{sec:hawkes-variants}).

\begin{modelcard}{Queue-reactive model}
\cardrow{What you calibrate}{The rate functions $r_L(x)$, $r_C(x)$, $r_M(x)$ for each of the first few
levels (and, in richer versions, their dependence on neighbouring queues); the price-move
probabilities $\pi_{\text{move}}$ and $\pi_{\text{reinit}}$.}
\cardrow{How}{Count events by queue size and divide by the time spent at that size; calibrate $\pi_{\text{move}}$
and $\pi_{\text{reinit}}$ to the observed volatility and to the continuation--reversal ratio.}
\cardrow{What you do with it}{Simulate a large-tick book that responds to one's own orders; compute fill
probabilities; compare placement tactics for a large order.}
\cardrow{How to picture it}{A toll booth whose customers and attendant react to the line: when it is
long, more drivers give up and leave; when it is short, cars are served faster.}
\end{modelcard}

\subsection{Decision rules from first-passage times}

Depletion times translate directly into execution decisions.
\begin{description}[leftmargin=1.5em,style=nextline,itemsep=2pt]
\item[Limit or market] If the expected time for a passive order to reach the front and be filled is
  longer than the horizon over which the trader's signal is valid, waiting forfeits the signal;
  crossing the spread is then the better choice despite its cost.
\item[When to cancel] If the estimated depletion time of the side on which an order rests collapses,
  or the race of Section~\ref{sec:cdl} turns against it, the order is about to be filled just
  before an adverse move, and should be cancelled or repriced.
\item[What a queue position is worth] \citet{moallemi2017} split the value of a resting order into a
  \emph{static} part---the spread earned at execution minus the adverse-selection cost, which grows
  with queue position---and a \emph{dynamic} part---the option value of holding a place in the queue
  that advances over time. They find that the front--back difference can be comparable to the
  bid--ask spread for some large-tick US stocks (for example BAC and CSCO), and smaller for others.
\item[Fill probability as survival] The time to fill is a waiting time that, for orders cancelled
  first, is never observed; survival models~\citep{cox1972}, which handle such \emph{censored}
  observations, are the natural statistical tool (Section~\ref{sec:deep-fill}).
\end{description}

Queue-depletion models thus serve execution in two ways: they estimate \emph{when} a resting order
will be filled, and, because a depletion is a price move, they indicate \emph{how informative} the
fill is likely to be. Section~\ref{sec:adverse} develops the second point. Depletion and price
prediction are related but distinct: a depletion model answers \emph{when} a level empties, and says
which way the price goes next only when it is used for the race between the two sides.

% ================================================================
\section{Hawkes processes}
\label{sec:hawkes}
% ================================================================

\paragraph{Works in calm, fails in stress.} The queue models of Section~\ref{sec:queue-models} assume
that orders arrive as independent streams at rates fixed over the sample. In a calm book that is a
reasonable approximation: the Cont--Stoikov--Talreja up-move probabilities of Table~\ref{tab:cst-up}
track the empirical frequencies closely. It breaks down exactly when it matters most. After a large
trade or a news release, rates jump to several times their average for a few seconds, and a model
that uses the day's average rate underestimates how fast queues empty (Example~\ref{ex:fpt}). The
missing ingredient is that events cause events. This section introduces the standard model of that
mechanism.

\subsection{Why not a Poisson process?}

The simplest model of events arriving at random times is the \emph{Poisson
process}~\citep{daley2003}. Divide time into tiny slices and, independently in each slice, let an
event happen with a small fixed probability; in the limit, events arrive at a constant average rate
$\mu$ per second. The number $N(u)$ of events between time $0$ and time $u$---and likewise the
number in any other interval of length $u$---has the distribution
\begin{equation}
  \P\big(N(u) = k\big) = \mathrm{e}^{-\mu u}\, \frac{(\mu u)^k}{k!}, \qquad k = 0, 1, 2, \dots
\end{equation}
\inwords{the chance of seeing exactly $k$ events in $u$ seconds depends only on the average rate
and the length of the interval---not on when the interval starts, and not on anything that happened
before it. The process has no memory: an event now makes the next one neither more nor less
likely. The waiting times between events are independent and exponentially distributed, and the
count in an interval has variance equal to its mean.}

For order-book events these properties fail, in three ways that can be seen in any day of data.
\begin{enumerate}[leftmargin=1.8em,itemsep=2pt]
\item \textbf{Events cluster.} Activity comes in bursts. A large trade is followed by a flurry of
  new orders, cancellations and further trades as other participants react; a news release does the
  same on a larger scale. Between bursts the market can be quiet for long stretches. In a Poisson
  process an event says nothing about the next one, so bursts can only arise by chance, and large
  ones are vanishingly rare.
\item \textbf{Counts vary far more than Poisson allows.} Because of clustering, the number of events
  per second has a variance much larger than its mean, whereas a Poisson count has variance equal to
  its mean~\citep{bacry2015}.
\item \textbf{Event types trigger each other.} A buy market order makes further buy market orders
  more likely, because large orders are split into many small ones and because other traders
  follow; a trade that empties a level triggers cancellations and new limit orders around it. A
  set of independent Poisson processes, one per event type, cannot represent any of this.
\end{enumerate}
A fourth, milder departure is that average activity changes predictably through the day, typically
highest near the open and the close. A Poisson process with a time-varying rate (an
\emph{inhomogeneous} Poisson process) can absorb that pattern, but it still cannot make events
cause other events, which is the defining feature of order flow.

\begin{example}[How rare is a burst under Poisson?]
Suppose events arrive at an average rate of $\mu = 10$ per second. Under a Poisson model the
probability of 30 or more events in a single second is about $2.5 \times 10^{-7}$. Over a 6.5-hour
trading day of 23\,400 seconds, the expected number of such seconds is about $0.006$---roughly one
every 170 trading days. In real order flow, seconds with several times the average activity occur
routinely around large trades and announcements. A model that calls them once-a-year events will
badly underestimate how quickly queues can empty (Section~\ref{sec:queue-models}).
\end{example}

A compact way to remember the difference: a Poisson process is \emph{drizzle}---drops fall
independently at a steady rate, and one drop says nothing about the next. Order flow is a
\emph{storm}---long dry spells broken by downpours, in which each burst of rain makes more rain likely
for a while.

Appendix~\ref{app:pp} gives a self-contained introduction to these ideas: the Poisson, renewal and
Hawkes processes, the conditional intensity, the branching ratio, and the statistics that tell the
three apart on real data, with simulated examples.

The Hawkes process keeps what is convenient about the Poisson process---events at random times,
described by a rate---but lets the rate \emph{rise after each event} and then decay back. It is the
smallest change to the Poisson model that makes events cause events. When every kernel below is
zero, it reduces to independent Poisson processes with rates $\mu_i$.

\subsection{The multivariate model}

Let $N_i(t)$ count events of type $i \in \{1,\dots,K\}$. A linear multivariate Hawkes
process~\citep{hawkes1971} has conditional intensity
\begin{equation}
  \lambda_i(t) = \mu_i + \sum_{j=1}^{K} \int_0^t \phi_{ij}(t-t')\, dN_j(t'),
\end{equation}
with baseline rates $\mu_i$ and non-negative kernels $\phi_{ij}$. Here $\lambda_i(t)$ is the
\emph{intensity}: the expected number of type-$i$ events per second at time $t$. The term
$dN_j(t')$ is nonzero only at the instants $t'$ when a type-$j$ event occurred, so the integral is
simply a sum over past events, and the formula can be read as
$\lambda_i(t) = \mu_i + \sum_j \sum_{\text{past type-}j\text{ events at } t'} \phi_{ij}(t - t')$.
\inwords{the current rate of type-$i$ events is a constant background rate $\mu_i$ plus a ``bump''
for every past event; each bump fades as that event recedes into the past, with a shape set by the
kernel $\phi_{ij}$. Events make further events more likely for a while, the way an earthquake
raises the rate of aftershocks or a viral post triggers reposts.}
The branching matrix $\Gamma_{ij} = \int_0^\infty \phi_{ij}(u)\,du$, the total area under a kernel,
is the expected number of type-$i$ events directly triggered by one type-$j$ event. The process is stationary when the spectral radius
$\rho(\Gamma) < 1$, the largest absolute eigenvalue of $\Gamma$. Then, with
$\mu = (\mu_1, \dots, \mu_K)^\top$ and $I$ the $K \times K$ identity matrix, the stationary mean
intensity vector $\bar\lambda = (\bar\lambda_1, \dots, \bar\lambda_K)^\top$ is
\begin{equation}
  \bar\lambda = (I - \Gamma)^{-1}\mu .
\end{equation}
\inwords{expand the inverse as a geometric series,
$(I-\Gamma)^{-1}\mu = \mu + \Gamma\mu + \Gamma^2\mu + \cdots$. The long-run event rate is the
background events, plus the events they trigger, plus the events \emph{those} trigger, and so on
through every generation. The series converges---the cascade dies out---exactly when
$\rho(\Gamma) < 1$. Readers who know discrete-time linear systems will recognise the condition:
it is the stability condition for $\psi_{k+1} = \Gamma \psi_k + \mu$, a feedback loop whose gain must be
below one.}

\paragraph{A note on notation.} Much of the literature writes the exponential kernel as
$\alpha\mathrm{e}^{-\beta u}$ and the (scalar) branching ratio as $n$. In this survey $\beta$ is
reserved for price impact (Section~\ref{sec:ofi}) and $n$ for event indices, so the decay rate is
written $\omega$ and the branching ratio $\Gamma$ (a $1 \times 1$ branching matrix). With
$\phi(u) = \alpha\mathrm{e}^{-\omega u}$, the branching ratio is $\Gamma = \alpha/\omega$.
Hawkes models have been applied extensively to order flow and price formation~\citep{bacry2015}.

\begin{example}[Endogeneity in two dimensions]
\label{ex:hawkes}
Take buy and sell market orders with $\mu = (1, 1)$ events per second and
\[
  \Gamma = \begin{pmatrix} 0.6 & 0.2 \\ 0.2 & 0.6 \end{pmatrix}.
\]
The eigenvalues are $0.8$ and $0.4$, so $\rho(\Gamma) = 0.8 < 1$. Then
$I - \Gamma = \begin{pmatrix} 0.4 & -0.2 \\ -0.2 & 0.4 \end{pmatrix}$ with determinant $0.12$, and
\[
  \bar\lambda = \frac{1}{0.12}\begin{pmatrix} 0.4 & 0.2 \\ 0.2 & 0.4 \end{pmatrix}
  \begin{pmatrix} 1 \\ 1 \end{pmatrix} = \begin{pmatrix} 5 \\ 5 \end{pmatrix}.
\]
Only one event in five is exogenous. The model says a great deal about clustering and nothing, by
itself, about direction: the two sides are symmetric.
\end{example}

\subsection{Reading the branching matrix}

The branching matrix is the economically interesting output of a fitted Hawkes model. Take two
event types, buy and sell market orders:
\[
  \Gamma = \begin{pmatrix} \Gamma_{BB} & \Gamma_{BS} \\ \Gamma_{SB} & \Gamma_{SS} \end{pmatrix}.
\]
$\Gamma_{BB}$ is the average number of further buy market orders directly triggered by one buy
market order: buys beget buys, for example because a large order is being executed in pieces.
$\Gamma_{SB}$ is the number of sell orders triggered by a buy: a cross effect, for example traders
leaning against a move. A large diagonal relative to the off-diagonal means trends in order flow; a
large off-diagonal means reversal. Asymmetry ($\Gamma_{BS} \ne \Gamma_{SB}$) means one side reacts
more strongly to the other than vice versa. Linear Hawkes models allow only excitation; inhibition
(an event making others \emph{less} likely) requires non-linear extensions. The shape of each kernel
$\phi_{ij}$ says how long an influence lasts: an exponential kernel forgets quickly, a power-law
kernel remembers for a long time.

\paragraph{An analogy: acoustic feedback.} A microphone near a loudspeaker picks up each sound and
plays it again, slightly quieter. If the loop gain is below one, every echo is fainter than the last
and a clap dies away; the closer the gain is to one, the longer the ringing lasts. At a gain of one or
more the echoes grow and the system howls. The branching ratio is that loop gain for order flow: the
average number of events that one event triggers. At $0.8$, a single exogenous event produces on average
$1/(1-0.8) = 5$ events in total, as in Example~\ref{ex:hawkes}; at $0.95$ it produces $20$.

\begin{example}[A chain reaction in the book]
\label{ex:chain}
A stylised sequence at the best ask, all within a few milliseconds:
\begin{center}\small
\begin{tabular}{rl}
\toprule
time & event \\
\midrule
0.0\,ms & a buy market order takes 10 of the 14 lots at the best ask (exogenous) \\
0.3\,ms & two sellers at the best ask cancel, fearing a move up (triggered) \\
0.6\,ms & a buyer joins the best bid, leaning with the flow (triggered) \\
0.9\,ms & a second buy market order takes the remaining lots; the best ask ticks up (triggered) \\
1.2\,ms & new sellers post at the new best ask; buyers follow at the new bid (triggered) \\
\bottomrule
\end{tabular}
\end{center}
One outside event produced four more and a price move. In a Hawkes model each row raises the intensity
of the event types in later rows---entries $\Gamma_{ij}$ for cancellation after trade, trade after
trade, limit order after cancellation---and the fitted kernels say how quickly that influence fades.
\end{example}

\subsection{What Hawkes models are for}

Hawkes models represent clustering, self- and cross-excitation, short-term persistence and the
split between exogenous and endogenous activity. They are valuable for intensity forecasting and
event simulation even when directional alpha is weak. A model can reproduce the temporal
statistics of order flow very well while adding little information about the sign of the next
price move. Hawkes models should therefore be evaluated separately for intensity forecasting,
simulation, price prediction and execution prediction (Table~\ref{tab:three-problems}).

An execution use deserves emphasis. The intensity of aggressive flow on one side is a natural
input to a cancellation decision for a passive order on that side: a spike in sell-market-order
intensity raises the probability that a resting bid is about to be filled at a bad time. This is
an adverse-selection question (Section~\ref{sec:adverse}), not a direction question.

\subsection{From events to prices}

The simplest bridge from event counts to prices treats each mid move as an event. Let $N^{+}(t)$
and $N^{-}(t)$ count the upward and downward mid moves up to time $t$, each of size $\delta_m$. Then
\begin{equation}
  m_t = m_0 + \delta_m\,\big(N^{+}(t) - N^{-}(t)\big).
\end{equation}
\inwords{the price is a running tally of up-ticks minus down-ticks; modelling the two counting
processes models the price.}

Write $\lambda^{+}(t)$ and $\lambda^{-}(t)$ for the intensities of $N^{+}$ and $N^{-}$: the expected
numbers of up and down moves per second. \citet{bacry2013} make the two counts \emph{mutually}
exciting: an up move raises the intensity of down moves and vice versa,
\begin{equation}
  \lambda^{+}(t) = \mu_0 + \int_0^t \alpha\,\mathrm{e}^{-\omega (t-t')}\, dN^{-}(t'), \qquad
  \lambda^{-}(t) = \mu_0 + \int_0^t \alpha\,\mathrm{e}^{-\omega (t-t')}\, dN^{+}(t'),
\end{equation}
\inwords{each down move adds a bump of height $\alpha$ to the up-move rate, and vice versa; the bump
decays at rate $\omega$. The background rate on each side is $\mu_0$.}
The cross-kernel has total weight $\Gamma_{\times} = \alpha/\omega < 1$. This one mechanism
produces the familiar microstructure effects. Measured over short windows, prices look very
volatile because every up-tick tends to be followed by a down-tick (bid--ask bounce); over long
windows that back-and-forth cancels. Writing $\sigma^2(w)$ for the variance of price changes per unit time
measured over windows of length $w$, in units of ticks squared,
\begin{equation}
  \sigma^2(w) = \Lambda\left(\frac{1}{(1+\Gamma_\times)^2} + \Big(1-\frac{1}{(1+\Gamma_\times)^2}\Big)\,\frac{1 - \mathrm{e}^{-(\alpha+\omega) w}}{(\alpha+\omega)\,w}\right),
  \qquad \Lambda = \frac{2\mu_0}{1-\Gamma_{\times}} .
\end{equation}
\inwords{$\Lambda$ is the total rate of price moves. For very short windows the measured variance is
$\Lambda$---every move counts. For long windows it falls to $\Lambda/(1+\Gamma_\times)^2 = 2\mu_0/\big((1-\Gamma_\times)
(1+\Gamma_\times)^2\big)$, because moves that are immediately reversed contribute nothing to the
long-run change. The curve between the two is the ``signature plot'' that practitioners use to choose
a sampling frequency for volatility estimates.} Here cross-excitation \emph{damps} long-run variance.
Self-excitation, by contrast, inflates the variance of event \emph{counts} relative to their mean
rate~\citep{bacry2014}; which effect dominates for prices depends on the structure of the kernel
matrix, not on a single number. The two-dimensional version of the model also reproduces the Epps
effect, the drop in measured correlation between two assets when it is computed over very short
windows~\citep{bacry2013}.

\citet{bacry2014} extend the construction to four kernels linking trades and price changes---trades
exciting trades, prices reverting, trades moving prices, and prices triggering trades---and recover
the concave, relaxing impact of large orders executed in pieces. The same machinery was used early to
study book resilience~\citep{large2007} and to make a zero-intelligence order-book simulator produce
realistic spreads~\citep{munitoke2011}.

\subsection{What fitted Hawkes models find}

\paragraph{Branching ratios.} Estimates of the branching ratio for financial event streams have
been high since the mid-2000s, but the exact value is contested. On E-mini S\&P~500 futures,
\citet{filimonovsornette2012} find it rising from about $0.3$ around 1998--2000 to $0.7$--$0.8$ after
2007, with spikes to $0.9$. \citet{hardimanbercotbouchaud2013}, using power-law kernels, find the
kernel integrating to approximately one---a market at criticality. \citet{filimonovsornette2015}
argue that estimates near one are biased upward by outliers, kernel regularisation, edge effects
and regime shifts---fitting a constant background rate to a period in which the true background
varies; their
re-estimation gives values that later fluctuate between about $0.8$ and $1.1$. On CME NQ market
data, an exponential-kernel fit to packet arrivals gives about $0.80$ (Appendix~\ref{app:pp}).
The defensible summary is that a large majority of events are endogenous, and that the precise
number depends on the kernel family and on how non-stationarity is handled.

\paragraph{Kernel shape and long memory.} Non-parametric estimation~\citep{bacry2016} finds kernels
that decay roughly as power laws with exponent slightly above one (in simulation the method recovers
power laws over six decades of time). In
an eight-dimensional model of DAX and Bund futures (price moves, trades, limit orders and
cancellations, each split by side), trades, limit orders and cancellations mainly excite events of
the same type and side, while price moves on one side mainly excite price moves on the other---mean
reversion. The persistence of order flow is a long-documented fact: the signs of market orders form a
long-memory process with a Hurst exponent of about $0.7$~\citep{lillo2004}, explained mainly by the
splitting of large orders into many small ones~\citep{lillo2005,bouchaud2009}. Power-law Hawkes
kernels are the point-process expression of the same fact.

\subsection{Near-criticality}

When the branching ratio approaches one, appropriately rescaled Hawkes processes converge to
continuous limits. With light-tailed kernels the limit is an integrated Cox--Ingersoll--Ross process,
and the price model of \citet{bacry2013} converges to a Heston stochastic-volatility
model~\citep{jaisson2015}. With heavy-tailed kernels---power laws with tail $\propto u^{-(1+\zeta)}$ with $\zeta \in
(1/2, 1)$---the limit is a \emph{rough} fractional process with Hurst parameter $\zeta -
1/2$~\citep{jaisson2016}, which the authors offer as a microstructural foundation for rough
volatility. Empirically, log-volatility does behave like a fractional process with a Hurst exponent of
order $0.1$---for example $0.125$ for DAX and $0.082$ for Bund futures~\citep{gatheral2018rough}. The lesson for LOB modelling is qualitative: strong endogenous excitation, together with
long memory, generates volatility that is itself random and persistent. It is not that markets
literally sit at a singularity, nor that a single scalar multiplies volatility.

\subsection{State-dependent Hawkes models}

The standard specification lets intensity depend on past events but not on the current book. A
state-dependent form~\citep{morariupatrichi2022} couples the two; one convenient parametrisation is
\begin{equation}
  \tilde\lambda_i(t) = \lambda_i(t)\, \exp\!\left(\vartheta_i^\top X_{t-}\right),
\end{equation}
where $\tilde\lambda_i(t)$ is the state-dependent intensity, $\lambda_i(t)$ the Hawkes intensity
defined above, $\vartheta_i$ a vector of coefficients,
and the state $X_{t-}$ may contain the spread, queue imbalance, depth, volatility or cross-asset variables.
\inwords{take the Hawkes rate and multiply it by a factor that depends on the current book. The
dot product $\vartheta_i^\top X_{t-}$ is a weighted score of the book's features, as in logistic
regression, and the exponential turns any score into a positive multiplier, so the rate can never
go negative.}
The gain is realism; the cost is more parameters and more work to fit them. The relevant question is whether state
dependence provides incremental out-of-sample economic value sufficient to justify that cost.

\begin{modelcard}{State-dependent Hawkes}
\cardrow{What you calibrate}{Everything in the linear Hawkes model, plus the coefficient vector
$\vartheta_i$ that scales each intensity by the current book state.}
\cardrow{How}{Maximum likelihood, as for the linear model, with the book state $X_{t-}$ evaluated just
before each event; more parameters, so more data and regularisation are needed.}
\cardrow{What you do with it}{Intensity forecasts that react to the current spread, imbalance and
depth, for simulation and for timing passive orders.}
\cardrow{How to picture it}{The same aftershock model, with a dial set by the state of the book that
turns every rate up or down.}
\end{modelcard}

\subsection{Fitting Hawkes models in practice}
\label{sec:hawkes-fit}

\paragraph{Maximum likelihood.} The standard method chooses the parameters that make the observed
event times most probable~\citep{ozaki1979}. For a $K$-type process observed on $[0, t_{\text{end}}]$ the
log-likelihood is a sum over types,
\begin{equation}
  \ln \mathrm{Lik} = \sum_{i=1}^{K} \Big( \sum_{n:\ \text{event } n \text{ of type } i} \ln \lambda_i(t_n)
  - \int_0^{t_{\text{end}}} \lambda_i(t')\, dt' \Big),
\end{equation}
\inwords{for every type, reward the model for a high intensity at the moments type-$i$ events
happened, and penalise it for the total number of type-$i$ events it expected.}
Because $\lambda_i$ depends only on $\mu_i$ and on the kernels $\phi_{i1}, \dots, \phi_{iK}$ in row
$i$ of the kernel matrix, the $i$-th term involves only those parameters. As long as no parameter is
shared across rows, the fit splits into $K$ independent problems, one per target type, each with
$1 + K$ kernels' worth of parameters~\citep{wu2019qrhawkes,bacry2015}. For exponential kernels each
term is computed in one pass with a running sum (Appendix~\ref{app:pp-mle}), so the cost grows only
linearly with the number of events~\citep{bacry2015}.

Three practical rules follow. First, parameters have natural bounds---$\mu_i > 0$, jumps
$\alpha_{ij} \ge 0$, decay $\omega > 0$---and the optimiser should enforce them. Second, nothing in the
likelihood forces the fitted $\Gamma$ to satisfy $\rho(\Gamma) < 1$; it must be checked after fitting,
and a fit at or above one read as a sign of non-stationarity or a wrong kernel rather than as a fact
about the market (Appendix~\ref{app:pp-branching}). Third, the likelihood is not convex in the decay
rates, but if the decays are fixed in advance---for example a few exponentials with chosen time
scales---it becomes convex in the background rates and jumps~\citep{wu2019qrhawkes}, so the remaining
fit has a single optimum.

\paragraph{Expectation--maximisation.} Direct maximisation can be unstable because the likelihood is
flat in some directions. The branching structure suggests a different route~\citep{veen2008em}: if we
knew which earlier event triggered each event, or that it was a background event, the parameters would
be simple averages. The EM algorithm alternates between guessing those parents probabilistically and
re-estimating. For one type with kernel $\alpha\mathrm{e}^{-\omega u}$~\citep{lewis2011em}:
\begin{itemize}[leftmargin=1.5em,itemsep=1pt]
\item \emph{E step:} the probability that event $n$ was triggered by earlier event $n'$ is
  $\alpha\mathrm{e}^{-\omega (t_n - t_{n'})}/\lambda(t_n)$, and the probability that it was a background
  event is $\mu/\lambda(t_n)$;
\item \emph{M step:} $\mu$ is the expected number of background events divided by $t_{\text{end}}$; $\Gamma$ is the
  expected number of triggered events divided by the number of events; and $1/\omega$ is the
  probability-weighted average of the parent--child delays $t_n - t_{n'}$; then $\alpha = \Gamma\omega$.
\end{itemize}
\inwords{give each event a fractional family tree, then count: how many events had no parent, how many
children an event has on average, and how long children take to arrive.}

\paragraph{Moments and cumulants.} A second family avoids the likelihood altogether and matches
statistics of the counts. The Wiener--Hopf method estimates kernel shapes from the empirical
second-order statistics of the event streams~\citep{bacrymuzy2016wh} (it underlies the non-parametric
results of~\citet{bacry2016}); the cumulant method NPHC estimates only the integrated kernels---the
branching matrix itself---from second- and third-order integrated cumulants, with no assumption on
kernel shape~\citep{achab2018cumulants}; and method-of-moments calibration with closed-form moments
can be rolled forward through time cheaply~\citep{dafonseca2014}.

\paragraph{The branching ratio from count dispersion.} The simplest moment estimator needs only
counts. Split the data into windows of length $w$, and compute the mean and variance of the number of
events per window; their ratio is the Fano factor $D(w)$ of Appendix~\ref{app:pp-fano}. For a
one-type Hawkes process whose correlations die out on time scales much shorter than $w$,
\citet{hardiman2014branching} show that
\begin{equation}
  \Gamma \approx 1 - \sqrt{1/D(w)} .
\end{equation}
\inwords{self-excitation makes counts burstier than Poisson; the more bursty they are, the closer the
branching ratio is to one. A Poisson stream ($D = 1$) gives $\Gamma = 0$.}
The estimate is biased low for finite windows and exact only as $w$ grows without bound, and because
the sample variance can be zero the authors report medians. Applied to the Fano factors of CME NQ
packet data reported in Appendix~\ref{app:pp-fano}---about 52 at $w = 0.1$\,s and about a thousand at
$w = 60$\,s---the formula gives $\Gamma \approx 0.86$ and $\Gamma \approx 0.97$, against $0.80$ from
the exponential-kernel likelihood fit on the same kind of data. Two cautions make this a teaching
example rather than a measurement. The estimate keeps rising with $w$, the same symptom that a slowly
varying background rate produces~\citep{filimonovsornette2015}. And the same stream with its gaps
shuffled---which destroys all self-excitation but keeps the gap distribution---has $D \approx 17$, for
which the formula would still report $\Gamma \approx 0.76$: dispersion alone does not prove
excitation, because a renewal process with very irregular gaps is also over-dispersed.

\paragraph{Non-stationarity.} A single fit over a long window assumes a constant background rate.
Because activity follows intraday patterns and changes across regimes, such fits overstate
endogeneity~\citep{filimonovsornette2015}. In practice models are fitted over short windows, with a
time-varying background, or with the background modelled explicitly.

\paragraph{Online use.} With exponential kernels the intensity is updated in constant time per event,
which makes Hawkes intensities usable inside a low-latency trading loop. The parameters change more
slowly. One workable scheme, which we describe rather than cite, holds the decay rates fixed during a
session and re-fits the background rates and jumps on a rolling window, starting each fit from the
previous solution; with fixed decays each re-fit is a convex problem. Fully online estimators also
exist: for example, \citet{yang2017online} update non-parametric kernels with projected gradient steps
on a discretised, regularised likelihood, with regret growing only logarithmically in time. Whatever
the scheme, the gradient of the log-likelihood contains the compensator term $-\int \lambda_i$ as well
as the event term, so an update that looks only at the events (``the rate was high when events
happened'') and ignores how many events the model expected will push the intensities up without limit.

\paragraph{Goodness of fit: time rescaling.} A fitted point process can be checked with a classical
transformation. Replace each event time $t_n$ by the expected number of events up to it under the
fitted model, $\int_0^{t_n} \widehat\lambda(t')\,dt'$. If the model is right, the transformed times form
a Poisson process with rate one, so the gaps between them are independent and exponential with mean
one~\citep{ogata1988,brown2002}. A quantile--quantile plot of those gaps against the exponential
distribution, or a Kolmogorov--Smirnov test, shows where the model fails---for example, too few very
short gaps means the model underestimates clustering.

\paragraph{Simulation.} Hawkes paths are generated by \emph{thinning}~\citep{ogata1981}: propose a
candidate event using an upper bound on the current intensity, which is valid because exponential
intensities only decay between events, and accept it with probability equal to the true intensity
divided by the bound. Simulated paths are how a fitted model answers questions that have no closed
form, such as the probability that a queue empties within a given latency.

\paragraph{Software.} The open-source library tick~\citep{bacry2018tick} implements maximum-likelihood
fits with exponential and sum-of-exponential kernels, EM, Wiener--Hopf and cumulant estimators, and
simulation.

\begin{modelcard}{Linear (multivariate) Hawkes}
\cardrow{What you calibrate}{Background rates $\mu_i$; the kernels $\phi_{ij}$---for exponential kernels,
the jumps $\alpha_{ij}$ and decay $\omega$---and hence the branching matrix $\Gamma$.}
\cardrow{How}{Maximum likelihood, split into one problem per target type, with bounded parameters and
a one-pass running sum for exponential kernels; or EM, Wiener--Hopf, cumulant and dispersion estimators.
Fit on short or rolling windows, check $\rho(\Gamma) < 1$ afterwards, and test the fit by time
rescaling.}
\cardrow{What you do with it}{Forecast event intensities and, from them, drift, expected OFI and
depletion times; simulate order flow by thinning; monitor how much activity is endogenous; raise a
cancellation alarm when the aggressive intensity against a resting order spikes.}
\cardrow{How to picture it}{Aftershocks after an earthquake, or acoustic feedback: each event raises
the rate of further events, which fades; the branching ratio is the loop gain.}
\end{modelcard}

\subsection{From intensities to decisions}
\label{sec:hawkes-decisions}

A fitted Hawkes model gives, at every instant, the expected rate of each kind of event. Four
quantities a trader can use follow directly; each is valid only over horizons short enough that the
intensities do not change much, because they decay between events.

\paragraph{Drift.} In the price model of the previous subsection, up and down moves of size
$\delta_m$ arrive at rates $\lambda^+(t)$ and $\lambda^-(t)$, so over a short horizon $H$ the expected
mid change is approximately
\[
  \E[m_{t+H} - m_t] \approx \delta_m\,\big(\lambda^+(t) - \lambda^-(t)\big)\, H .
\]
\inwords{if up moves are currently arriving faster than down moves, the price is expected to drift up,
by one tick per surplus up move.}

\paragraph{Expected order-flow imbalance.} Model six event types at the best quotes: limit orders,
cancellations and market orders on each side, with intensities $\lambda^{L,b}$, $\lambda^{C,b}$,
$\lambda^{M,b}$ (bid queue added to; bid queue cancelled; market \emph{buy}, which takes from the ask)
and $\lambda^{L,a}$, $\lambda^{C,a}$, $\lambda^{M,a}$ (ask queue added to; ask queue cancelled; market
\emph{sell}, which takes from the bid). With $\bar v$ the average event size and the best prices
unchanged, the expected OFI per second is
\[
  \bar v\,\Big( \underbrace{\lambda^{L,b} - \lambda^{C,b} - \lambda^{M,a}}_{\text{net change of the bid queue}}
  \;-\; \underbrace{\big(\lambda^{L,a} - \lambda^{C,a} - \lambda^{M,b}\big)}_{\text{net change of the ask queue}} \Big).
\]
\inwords{the same score as Section~\ref{sec:ofi}, but computed from expected rates instead of observed
events: additions to the bid and removals from the ask count as buying pressure.} It is a forecast of
OFI, which by the regression of Section~\ref{sec:ofi} is a forecast of the price change.

\paragraph{Depletion time.} The bucket picture of Section~\ref{sec:queue-models} needs an inflow and
an outflow. For the best ask, $r_{\text{in}} = \bar v\,\lambda^{L,a}$ and $r_{\text{out}} = \bar
v\,(\lambda^{M,b} + \lambda^{C,a})$, so with the current intensities
\[
  \E[\tau^a] \approx \frac{v^a_{1,t}}{r_{\text{out}} - r_{\text{in}}}, \qquad r_{\text{out}} > r_{\text{in}} .
\]
When this falls below the trader's own latency $T_{\text{lat}}$, a sell order resting in that ask queue
is in danger: if the trader decides to cancel, the level may be swept---filling the order just before
the price moves up---before the cancellation arrives. The
full probability $\P(\tau^a \le T_{\text{lat}})$ has no closed form under a Hawkes model, but can be
estimated by simulating from the current state by thinning.

\paragraph{A branching-ratio monitor.} Re-fitting $\Gamma$ on a rolling window gives a running
measure of how much of current activity is self-generated. A rising value signals a market in which
events are feeding on each other. Two cautions apply. A rise can be produced by a changing background
rate rather than by more excitation~\citep{filimonovsornette2015}; and published estimates differ
widely by instrument, period and method (Section~\ref{sec:hawkes}), so a fixed universal alarm level is
not justified. A threshold should be set from the instrument's own history, and the monitor is a
regime indicator, not a forecast of direction.

\subsection{Variants of the Hawkes model}
\label{sec:hawkes-variants}

The linear model with exponential kernels is a starting point. Each variant below relaxes one of its
assumptions; Table~\ref{tab:hawkes-variants} gives the model card for each.

\paragraph{Sums of exponentials.} Power-law kernels capture long memory but lose the constant-time
update. A power law can be approximated over many decades by a small number of exponentials with
geometrically spaced decay rates---for an exponent of two, about 1.7 exponentials per decade
\citep{bochud2007}---and each exponential keeps its own running sum, so the intensity is still updated
in constant time per event. With the decay rates fixed, the likelihood is convex in the remaining
parameters~\citep{wu2019qrhawkes}.

\paragraph{Nonlinear and inhibitory kernels.} Linear Hawkes kernels can only excite. A nonlinear Hawkes
process passes the sum through a function that keeps the intensity non-negative,
$\lambda_i(t) = h\big(\mu_i + \sum_j \int \phi_{ij}(t-t')\, dN_j(t')\big)$, which allows negative
kernels, that is, events that make others \emph{less} likely; general conditions for a stationary
version were given by~\citet{bremaud1996}. In limit-order-book data, \citet{lu2018} observe inhibition
effects that lead them to nonlinear Hawkes models, and \citet{rambaldi2017marked} find that limit orders
(and cancellations) on one side inhibit the same type on the other side, which they read as a
fair-price effect.

\paragraph{Marks: the size of the event.} A 100-lot trade should excite more than a 1-lot trade, but
the obvious fix---multiplying the kernel by a function of size---is rejected by the data: on Bund and
DAX futures the kernels for different trade sizes have different \emph{shapes}~\citep{rambaldi2017marked}.
The authors instead split events into volume classes and treat each class as its own event type. They
find that large trades excite limit orders on the opposite side, and that large orders are more
exogenous than small ones.

\paragraph{Regimes.} In a Markov-modulated Hawkes process a hidden state switches among a few regimes,
each with its own background rate and decay, and the parameters are fitted by EM~\citep{wang2012mmhp}.
It was developed for earthquake sequences (main shocks, aftershocks, quiescence); the analogue for markets
is a model whose quiet and busy regimes have different parameters, rather than one set of parameters
averaged over both.

\paragraph{Quadratic Hawkes and the Zumbach effect.} Volatility has a time asymmetry: past volatility
measured over long horizons predicts future short-horizon volatility better than the reverse---the
\emph{Zumbach effect}~\citep{zumbach2009}. A Hawkes intensity that depends only on past event counts
cannot produce it. The quadratic Hawkes model of~\citet{blanc2017qhawkes} adds a term in the
\emph{square} of a moving average of past signed price changes, so that a sequence of moves in the same
direction raises activity more than the same number of moves that cancel out. Calibrated on five-minute
bins of 133 NYSE stocks over 2000--2009, it reproduces the time asymmetry, and it generates long memory
in volatility without requiring the process to be critical.

\paragraph{Queue-reactive Hawkes.} \citet{wu2019qrhawkes} combine the two ideas of this part: the
background rates depend on the current queue size, as in Section~\ref{sec:qr}, and events excite each
other through sum-of-exponential kernels. On Bund and DAX futures the Hawkes term greatly improves the
pure queue-reactive model, both in the distribution of times between events and in the distribution of
queue sizes; market-order and price-change rates are mainly functions of the volume imbalance, while
limit and cancellation rates are not.

\paragraph{Cross-asset Hawkes.} Events in one instrument can excite events in another; this is the
point-process form of lead--lag, discussed in Section~\ref{sec:cross}.

\begin{table}[ht]
\centering\footnotesize
\caption{Model cards for the Hawkes variants.}
\label{tab:hawkes-variants}
\setlength{\tabcolsep}{3pt}
\begin{tabularx}{\linewidth}{>{\raggedright\arraybackslash}p{0.14\linewidth}*{4}{>{\raggedright\arraybackslash}X}}
\toprule
Variant & What you calibrate & How & What you do with it & How to picture it \\
\midrule
Sum of exponentials & jumps per (pair of types, decay); decays fixed on a geometric grid & convex likelihood once decays are fixed & long memory with constant-time updates, in software or hardware & a power law drawn with a few rulers of different lengths \\
Nonlinear / inhibitory & signed kernels; the link function & maximum likelihood (poorly convex) & represent events that discourage others & a thermostat as well as an amplifier \\
Marked (volume classes) & one set of kernels per size class & as for a multivariate model, one type per class & let large trades matter more, with their own time profile & a big stone makes bigger, longer ripples \\
Markov-modulated & per-regime background and decay; switching rates & EM over hidden regimes & separate quiet and busy regimes & weather that switches between drizzle and storm \\
Quadratic (Zumbach) & self-excitation kernel; return-feedback kernel & on binned returns (GMM, then likelihood) & volatility forecasts with the right time asymmetry & a trend, not just activity, wakes the market up \\
Queue-reactive Hawkes & size-dependent background; excitation kernels & maximum likelihood with fixed decays & simulate queues with realistic clustering & the toll booth of Section~\ref{sec:qr}, with rush hours \\
\bottomrule
\end{tabularx}
\end{table}

\subsection{Neural point processes}
\label{sec:npp}

\paragraph{Where we are.} Everything so far in this section has the same structure: the rate of
type-$i$ events is a background rate plus a sum of bumps, one per past event, each with a fixed shape
$\phi_{ij}$. The variants of Section~\ref{sec:hawkes-variants} change one ingredient at a time---the
shape of the bump, its sign, its dependence on size or on the queue. A neural point process changes
all of them at once: it keeps the idea of an intensity, but lets a neural network compute it from the
history. This subsection explains what that means, step by step.

\paragraph{What the Hawkes formula cannot say.} Three limitations of the linear Hawkes model motivate
the step.
\begin{enumerate}[leftmargin=1.8em,itemsep=2pt]
\item \textbf{Effects only add.} Each past event contributes its own bump, independently of the
  others. In the chain reaction of Example~\ref{ex:chain}, a trade followed by two cancellations at the
  same level is far more alarming than either alone: together they say the level is about to go.
  A sum of separate bumps cannot express ``A \emph{and then} B matters more than A plus B''.
\item \textbf{The shape is fixed in advance.} The analyst chooses exponential, power law, or a sum of
  exponentials, and the data only set the parameters.
\item \textbf{The rate can only jump up and then decay.} Between events, a linear Hawkes intensity
  with positive kernels can only fall. Real flow can do the opposite: after a sweep, new limit orders
  may arrive slowly at first and then faster as participants re-quote.
\end{enumerate}

\paragraph{The recipe.} Every neural point process has three parts.
\begin{description}[leftmargin=1.5em,style=nextline,itemsep=2pt]
\item[1. Describe each event as a vector] Event $n$ is turned into a vector of numbers (an
  embedding) built from its type, the time since the previous event $t_n - t_{n-1}$, and any marks
  such as size or side.
\item[2. Summarise the history] A sequence model reads the embeddings of events $1, \dots, n$ in order
  and produces a \emph{history vector} $\mathbf{z}^{\text{hist}}(t)$ that summarises everything that
  happened before time $t$. This is the role played in a Hawkes model by the running sums of decaying
  bumps.
\item[3. Turn the summary into intensities] For each event type $i$, the intensity is
\end{description}
\begin{equation}
  \lambda_i(t) = \operatorname{softplus}\!\big(\langle \mathbf{o}_i, \mathbf{z}^{\text{hist}}(t)\rangle\big),
  \qquad \operatorname{softplus}(y) = \ln\big(1 + \mathrm{e}^{y}\big),
\end{equation}
where $\mathbf{o}_i$ is a vector of learned output weights for type $i$ and $y$ stands for any real
number.
\inwords{score the current history with a weighted sum, one set of weights per event type, and pass
the score through the softplus, a smooth function that is always positive and close to the score
itself when the score is large. Rates must be positive; softplus guarantees it, as the exponential
did in the state-dependent model.}
A linear Hawkes model is the special case in which the ``summary'' is the vector of running sums
$\sum \phi_{ij}(t - t_{n'})$ and the output step simply adds them to $\mu_i$. A neural model replaces
the hand-written summary with a learned one.

\paragraph{Training: the same likelihood as Hawkes.} The parameters---the network weights and the
output weights $\mathbf{o}_i$---are fitted by maximising exactly the log-likelihood of
Section~\ref{sec:hawkes-fit}: reward high intensity at the times events happened, penalise the
total expected number of events $\int \lambda_i$. Gradient descent replaces the specialised
optimisers, because the network has many parameters. One new difficulty appears. For an exponential
kernel the integral $\int \lambda_i$ has a closed form; for a neural intensity it usually does not, and
is estimated numerically---commonly by evaluating the intensity at randomly sampled times between
events and averaging (Monte Carlo integration). Two lines of work remove this approximation:
\citet{omi2019} model the \emph{integral} of the intensity with a network constrained to be increasing,
and obtain the intensity as its derivative, so the likelihood is exact; \citet{shchur2020ifl} skip the
intensity altogether and model the probability distribution of the time to the next event directly,
for example as a mixture of simple distributions, which also makes sampling the next event exact.

\paragraph{Three ways to summarise the history.} The models in the literature differ mainly in step 2.
\begin{description}[leftmargin=1.5em,style=nextline,itemsep=2pt]
\item[Recurrent networks] An early and widely cited model, RMTPP~\citep{du2016rmtpp}, updates a recurrent
  network's state at each event and lets the intensity change between events through a fixed function
  of the elapsed time. The \emph{neural Hawkes process}~\citep{mei2017} makes the recurrent state itself
  evolve in continuous time: each memory cell of an LSTM decays exponentially between events towards a
  target value set at the last event. Because the target can be above or below the current value, the
  intensity can rise or fall between events, and one event can make another \emph{less} likely---the
  inhibition that linear Hawkes cannot represent.
\item[Attention] The Transformer Hawkes process~\citep{zuo2020thp} and the self-attentive Hawkes
  process~\citep{zhang2020sahp} let each event attend to all earlier events, as in the
  queries--keys--values box of Section~\ref{sec:deep}, with the time gaps encoded into the
  embeddings. The attention weights play the part of the kernel: they say how much each past event
  matters now, but are computed from the content of the events rather than fixed in advance.
  \citeauthor{zhang2020sahp} argue that the attention weights also give an interpretable measure of
  influence between event types.
\item[State-space models] The Mamba Hawkes process~\citep{gao2024mhp} summarises the history with a
  selective state-space model (Section~\ref{sec:deep}), which processes long sequences in time linear in
  their length.
\end{description}
\citet{shchur2021} survey the family and its design choices.

\begin{example}[What a learned summary can add]
Take the chain reaction of Example~\ref{ex:chain}. A two-type Hawkes model of trades and ask
cancellations gives the intensity of a further buy market order as background plus one bump for the
trade and one for each cancellation, whatever their order. A neural model reads the sequence ``buy
trade, cancel, cancel'' as a whole. If, in the training data, that particular pattern was usually
followed within a millisecond by a sweep of the level, the model can raise the intensity of buy market
orders far more than the sum of the separate bumps would, and leave it unchanged after the same three
events in a different order. Whether it learns this depends entirely on whether the pattern is
frequent enough in the data; the flexibility is useless without the examples.
\end{example}

\paragraph{Applications to order books.} The Transformer Hawkes process and the Mamba Hawkes process
both report results on a data set of financial transactions (buy and sell events of a stock over a
day), among other benchmarks~\citep{zuo2020thp,gao2024mhp}. Models built specifically for the book go
further. \citet{shi2022sdpnhp} compute a separate intensity for each order-book event type with
parallel continuous-time LSTMs and add an interaction between events and the state of the book; on
five days of LOBSTER data for three US stocks it outperforms both classical and neural point-process
baselines in likelihood, timing and event-type accuracy, and the authors note that their simulator
settings remain simple. \citet{lalor2025nhp} drive a twelve-event-type order-book simulator with an
LSTM-based neural Hawkes process and use it as the training environment for a reinforcement-learning
market maker.

\paragraph{What is lost.} The flexibility has three costs, and each matters in this survey's setting.
\begin{itemize}[leftmargin=1.5em,itemsep=2pt]
\item \textbf{Interpretability.} A linear Hawkes fit returns a branching matrix whose entries have a
  meaning (Section~\ref{sec:hawkes}) and a branching ratio that can be monitored. A neural model returns
  weights; quantities such as the fraction of endogenous events have no direct equivalent.
\item \textbf{Speed.} An exponential Hawkes intensity is updated with a few multiplications per event.
  A neural intensity needs a forward pass of a network per event, and more if the intensity is needed
  between events---a cost that matters inside a trading loop (Section~\ref{sec:hardware}).
\item \textbf{Guarantees.} The stability condition $\rho(\Gamma) < 1$ has no simple analogue, so a
  simulated neural model can drift into regimes never seen in training.
\end{itemize}
The evaluation follows Principle~3. The baseline is a multivariate Hawkes model, ideally with
sum-of-exponential kernels, fitted to the same events; the comparison is the held-out log-likelihood per
event, the time-rescaling test of Section~\ref{sec:hawkes-fit} (which applies to any model that outputs
an intensity), and, for a trading use, the economic value of the forecasts. A neural point process
earns its place when it improves these by more than its cost in speed and interpretability.

\begin{modelcard}{Neural point processes}
\cardrow{What you calibrate}{The embedding of events, the weights of the sequence model that summarises
the history (recurrent, continuous-time recurrent, attention or state-space), and the output weights
$\mathbf{o}_i$ that turn the summary into an intensity for each event type.}
\cardrow{How}{Gradient descent on the point-process log-likelihood, with the integral of the intensity
estimated by Monte Carlo, or made exact by modelling the integral or the next-gap distribution
directly. Compare on held-out likelihood and by time rescaling against a multivariate Hawkes baseline.}
\cardrow{What you do with it}{Predict the type and time of the next event; simulate order flow,
including as a training environment for execution or market-making agents; capture interactions
between events that additive kernels miss.}
\cardrow{How to picture it}{A Hawkes model whose bumps are drawn by a network that has read the whole
history in order, instead of by a fixed formula: it can notice that ``trade, cancel, cancel'' means
more than the three events separately.}
\end{modelcard}

% ================================================================
\section{Deep learning for limit order books}
\label{sec:deep}
% ================================================================

Deep learning entered LOB research around 2017 and now dominates the volume of published work. This
section describes what the main families of models actually do, what the large benchmark studies
have found when the models are re-tested, and which design choices matter. Readers new to the
architectures will find short definitions of convolutional, recurrent and attention-based networks
in the glossary (Appendix~\ref{app:glossary}) and textbook treatments in~\citet{goodfellow2016}.

\subsection{Problem formulations}

Most deep LOB papers solve one of six problems:
\begin{description}[leftmargin=1.5em,style=nextline,itemsep=2pt]
\item[Direction classification] predict whether the mid will be up, down or stationary after a horizon $H$
  measured in events or seconds. Labels are usually built by comparing a smoothed future mid with the current one
  against a threshold; the choice of smoothing and threshold changes the problem, and a recent paper
  shows that coupling the smoothing window to the horizon biases the labels~\citep{berti2025tlob}.
\item[Multi-horizon forecasting] predict returns, or their direction, over several horizons at
  once---an ``alpha term structure''~\citep{kolm2023,zhangzohren2021}.
\item[Next-event prediction] predict the type, side, price and size of the next message---the
  natural task for message-level and point-process models~\citep{linna2025lobert}.
\item[Full-book forecasting] predict the whole future book, prices and volumes at several
  levels~\citep{arxiv240902277}.
\item[Fill-time estimation] predict the distribution of the time until a resting order is filled
  (Section~\ref{sec:deep-fill}).
\item[Generation] produce synthetic message streams or book snapshots (Section~\ref{sec:deep-gen}).
\end{description}

\paragraph{The benchmark data set.} A large share of the literature reports results on FI-2010
\citep{ntakaris2018}: five stocks from the Helsinki exchange over ten trading days in June 2010,
about 4.5 million messages, ten levels, already normalised and downsampled into blocks of ten
events. It made comparison possible, and it also made overfitting easy. The DeepLOB authors
themselves describe it as ``far too short, downsampled and taken from a less liquid
market''~\citep{zhang2019}. Results on it should be read with the comparability fields of
Section~\ref{sec:comparable} in mind.

\subsection{Convolutional and recurrent models}

\paragraph{First models.} \citet{tsantekidis2017} applied a convolutional network to sequences of
book snapshots, treating the book like an image whose two axes are time and price level. A
companion paper used an LSTM with 40 hidden units on the last 100 snapshots, reporting F1 scores
of about 61--66\% on FI-2010 against 48--56\% for a multilayer perceptron~\citep{tsantekidis2017lstm}.
\citet{dixon2018} applied a recurrent network to E-mini S\&P~500 Level-2 data to classify the next
price flip.

\paragraph{DeepLOB.} \citet{zhang2019} combined the two ideas and became the most widely used
reference architecture. Its input is the last 100 snapshots of ten levels (40 numbers each:
price and size for bid and ask at each level). Small convolutions first combine price with size at
each level, then bid with ask, then levels with each other; an Inception module applies parallel
convolutions over several time scales; a 64-unit LSTM summarises the sequence; and a softmax outputs
probabilities for down, stationary and up. On a year of London Stock Exchange data for five stocks it
reports F1 of about 70\%, 63\% and 61\% at horizons of 20, 50 and 100 events, with similar numbers on
five stocks not seen in training---evidence that it learns features that transfer. Its trading
simulation, as the authors state, uses mid-prices without transaction costs and is ``not \ldots a fully
developed, stand-alone trading strategy''.

\paragraph{Learning across many stocks.} \citet{sirignano2019dl} introduced a ``spatial'' network that
exploits the ordering of price levels to use information deep in the book with few parameters,
trained on 489 US stocks over twenty months on a GPU cluster; it outperforms logistic regression and a
standard network, especially in the tails of the distribution. \citet{sirignano2019} train a single
model on many stocks and find that it outperforms stock-specific models, which they interpret as
evidence of universal features of price formation.

\paragraph{Lighter models and normalisation.} Not every gain needs more capacity. The temporal
attention-augmented bilinear network (TABL) acts separately on the feature and time axes of the input
matrix and reaches competitive FI-2010 results with about $10^4$ parameters~\citep{tran2019tabl}.
Deep adaptive input normalisation (DAIN) learns how to normalise each input
feature~\citep{passalis2020dain}; on FI-2010 it raises the macro-F1 of a simple MLP from about 55\%
with standard z-scoring to about 68\%. Because the statistics of LOB data drift, how inputs are
normalised is a modelling decision in its own right.

\begin{modelcard}{Convolutional and recurrent models (DeepLOB family)}
\cardrow{What you calibrate}{The network weights (from about $10^4$ to $10^5$ parameters for TABL and
DeepLOB), and before that the input normalisation, the label threshold and the horizon.}
\cardrow{How}{Supervised training with a cross-entropy loss on labelled windows of snapshots
(down/stationary/up), with chronological train, validation and test splits.}
\cardrow{What you do with it}{Short-horizon direction probabilities, used as a signal or as one input
to an execution decision.}
\cardrow{How to picture it}{The book as an image---time across, price levels down---scanned by small
filters that find local patterns, followed by a memory that reads the patterns in order.}
\end{modelcard}

\subsection{Attention and spatial-temporal models}

\paragraph{Attention.} An attention layer lets each element of a sequence form a weighted average of
all the others, with weights computed from the data~\citep{vaswani2017}. Applied to LOBs, the
sequence can run along time, along price levels, or along assets. A model that attends along more
than one of these axes is called \emph{spatial-temporal}: ``space'' is the price ladder (and, in
multi-asset models, the set of instruments), ``time'' is the sequence of snapshots or events.

\begin{tcolorbox}[breakable, enhanced, colback=black!3, colframe=black!45, boxrule=0.5pt, arc=1.5pt,
  fonttitle=\bfseries\small, title={Queries, keys and values on an order book}, fontupper=\small]
Each event $n$ is first turned into a vector of numbers (its embedding). Three learned matrices turn
that vector into a \emph{query} $\mathbf{z}^{\text{qry}}_n$ (``what am I looking for?''), a \emph{key}
$\mathbf{z}^{\text{key}}_n$ (``what do I contain?'') and a \emph{value} $\mathbf{z}^{\text{val}}_n$
(``what do I pass on?''), each of dimension $\dim_{\text{emb}}$. Event $n$ attends to earlier event
$n'$ with weight
\[
  \mathrm{att}_{n n'} = \frac{\exp\big(\langle \mathbf{z}^{\text{qry}}_n, \mathbf{z}^{\text{key}}_{n'}\rangle/\sqrt{\dim_{\text{emb}}}\big)}
  {\sum_{n'' \le n} \exp\big(\langle \mathbf{z}^{\text{qry}}_n, \mathbf{z}^{\text{key}}_{n''}\rangle/\sqrt{\dim_{\text{emb}}}\big)},
  \qquad \text{output}_n = \sum_{n' \le n} \mathrm{att}_{n n'}\, \mathbf{z}^{\text{val}}_{n'} .
\]
\inwords{each event scores every earlier event by how well its key matches the query, turns the scores
into weights that sum to one (the softmax), and takes the weighted average of their values.}

\emph{A stylised sweep.} The current event is a buy market order that empties the best ask. If training
has taught the model that sweeps matter when sellers have been leaving, the sweep's query matches the
keys of the ask cancellations in the last few milliseconds, which receive most of the weight; the bid
limit orders and the quiet stretch before them receive little. The output is, in effect, a summary
``sellers were pulling before this sweep''---the chain reaction of Example~\ref{ex:chain}, read off by
a learned look-up rather than a fixed kernel.
\end{tcolorbox}

\paragraph{TransLOB.} \citet{wallbridge2020} combined five dilated causal convolutions with two
Transformer blocks using masked (causal) self-attention, and reported FI-2010 F1 of 81--92\% across
horizons. An independent re-implementation obtained 59\% on the same data~\citep{prata2024}---the
largest gap between claimed and reproduced performance among the fifteen models tested there, and a
reminder that results on one small data set are fragile.

\paragraph{Axial and dual attention.} Axial-LOB~\citep{kisiel2022axial} applies gated attention
separately along the time and feature axes, keeping the parameter count below $10^4$. TLOB
\citep{berti2025tlob} alternates temporal self-attention (across snapshots) with spatial
self-attention (across the features of a snapshot); its companion MLPLOB replaces attention by
multilayer perceptrons that mix along the two axes, and matches TLOB on several data sets. Three
observations from TLOB are instructive. Performance on Nasdaq and Bitcoin data falls steeply with the
horizon---for INTC, F1 drops from about 80\% at 10 events to about 50\% at 100---whereas on FI-2010 it
\emph{rises} with the horizon, a sign of the data set's peculiarities. Predictability on the same stock
declined between 2012 and 2015. And when the label threshold is set to the average spread, so that
only moves large enough to pay for crossing count, F1 drops to 31--41\% (TSLA, horizons of 50 to 200
events). The authors conclude that the
methods ``are not sufficiently mature for practical deployment''.

\paragraph{Patches, full books and many assets.} LiT~\citep{lit2025} cuts the book into structured
patches---one side of the book over a short time window---as vision Transformers do with images, and
applies it to Binance cryptocurrency data. \citet{arxiv240902277} forecast the entire five-level book
with a sequence-to-sequence Transformer and a loss that penalises forecasts violating the ordering of
prices; on one day of data a linear model is slightly better on the overall forecasting error and
on volumes, but worse on the mid-price and much worse at respecting the ordering of prices. OF-MATNet
\citep{ofmatnet} feeds multi-level order-flow imbalance from many Nasdaq assets into a Transformer
that attends along time, assets and levels, choosing peer assets by rolling Granger-causality tests,
and reports improvements in out-of-sample $R^2$ in over 90\% of cases on 110 assets. HLOB
\citep{briola2025hlob} builds a graph over volume levels and processes it with a homological
convolutional network.

\paragraph{Multi-horizon models and training hardware.} \citet{zhangzohren2021} attach a
sequence-to-sequence decoder with attention to the DeepLOB encoder to produce forecasts at several
horizons at once, and report that training on Graphcore intelligence processing units took about
15\% of the wall-clock time of a GPU.

\begin{modelcard}{Attention and spatial-temporal Transformers}
\cardrow{What you calibrate}{The query, key and value matrices of every attention layer, the feed-forward
weights, and the choice of axes (time, levels, assets) along which to attend.}
\cardrow{How}{Supervised training as for the convolutional models; reproduction studies show that the
result depends heavily on the data set and labelling, so re-testing on new data is part of
calibration.}
\cardrow{What you do with it}{Direction and multi-horizon forecasts, full-book forecasts, and
multi-asset models that choose which peers to read.}
\cardrow{How to picture it}{Every event asks a question of every earlier event and listens most to the
ones whose answers match.}
\end{modelcard}

\subsection{State-space models}

Structured state-space sequence models such as S4~\citep{gu2022s4} and Mamba~\citep{gu2024mamba}
process long sequences with a learned linear recurrence---the same mathematical object as the
classical state-space model of control theory~\citep{kalman1960}---computed efficiently in parallel.
They are attractive for event streams, which are long and irregular. Apart from a hardware-aware design for cryptocurrency tick data~\citep{popov2026ssm}, we could not
identify a widely cited state-space model for LOB forecasting at the time of writing; this is an open
direction rather than an established result (Section~\ref{sec:agenda}).

\begin{modelcard}{State-space sequence models}
\cardrow{What you calibrate}{The matrices of a learned linear recurrence and, for selective models such
as Mamba, the functions that set its step size from the input.}
\cardrow{How}{Supervised training by gradient descent, with the recurrence computed in parallel during
training.}
\cardrow{What you do with it}{Long-context forecasts on long, irregular event streams; for LOB data this
is an open direction rather than an established result.}
\cardrow{How to picture it}{A Kalman-style running state, updated at every event, whose rate of
forgetting is learned.}
\end{modelcard}

\subsection{What the large benchmark studies show}
\label{sec:benchmarks-dl}

Individual model papers report results on their own terms. Two studies re-test many models under
common conditions, and two more vary the input or the evaluation while holding the model fixed; their
findings are more informative than any single paper.

\paragraph{LOBCAST.} \citet{prata2024} re-implemented fifteen published models and tested them on
FI-2010 and on new Nasdaq data from 2021 and 2022. Their headline is that ``all models exhibit a
significant performance drop when exposed to new data'': F1 on the new data lies between 48\% and 61\%
for every model. The best model on the new data, BiNCTABL, loses about a fifth of its FI-2010
score. Five of the six best models use attention. A simple MLP ranks near the bottom but degrades the
least. Parameter counts range from about $10^4$ to $10^7$, with no clear relation between size and
degradation. A trading simulation shows that profitability ``is far from guaranteed''.

\paragraph{Simple models on the same inputs.} \citet{briola2020} compare a random classifier,
logistic regression, an MLP, LSTMs and a CNN--LSTM on identical inputs for one large-tick stock. The
MLP matches or exceeds the CNN--LSTM at all three horizons, while logistic regression is clearly worse.
Linear models are not generally competitive with deep ones on raw book inputs; small non-linear
models often are.

\paragraph{Representation.} \citet{lucchese2024} hold the DeepLOB architecture fixed and change only
the input, on ten Nasdaq stocks. Models fed order-flow inputs or a new volume-on-a-price-grid
representation ``considerably'' outperform those fed the raw level-by-level book, and predictability
is detected up to about 50 order-book events ahead at the 99\% confidence level, using model
confidence sets~\citep{hansen2011mcs} to compare models. The authors stress that mid-to-mid returns
``are not tradable in practice''. \citet{kolm2023}, on 115 Nasdaq stocks, find that networks
trained on order flow ``significantly outperform most models trained directly on order books'', and
that the effective horizon of stock-specific forecasts is about two average price changes. In a later
practitioner-oriented account, as reported in coverage of the authors' 2026 award, very deep models
``don't really help'' for this task, and careful time-stamping and stationary transformations of inputs
mattered more~\citep{kolm_westray}. \citet{wu2021robust} show that level-based inputs are fragile:
adding orders in empty price ticks drops DeepLOB's accuracy from about 77\% to about 51\%, and they
propose representations on a fixed price grid as an alternative.

\paragraph{From forecasts to trades.} \citet{briola2024} retrain DeepLOB on fifteen Nasdaq stocks
grouped by tick size and introduce a transaction-level score,
\begin{equation}
  \mathrm{TO} = \frac{N_{\text{correct}}}{N_{\text{possible}} + N_{\text{predicted}} - N_{\text{correct}}},
\end{equation}
where $N_{\text{possible}}$ counts the round-trip trades (open, then close) that the true price path
would have allowed, $N_{\text{predicted}}$ those the forecasts would have triggered, and
$N_{\text{correct}}$ those in both sets.
\inwords{$\mathrm{TO}$, which its authors call $p_T$, is the overlap between the trades the model would make and the trades that were
actually there to be made, as a fraction of their union---a score of 1 means the model traded exactly
the available opportunities, 0 means no overlap.} The values are sobering: between about 0.01 and 0.19
across stocks and horizons, and close to zero for small-tick stocks once a confidence threshold is
applied, even where F1 is high. Predictability, by both F1-type and transaction-level measures, is
higher for large-tick stocks---those whose average spread is below about 1.5 ticks---than for
small-tick stocks, whose spread exceeds about 3 ticks. The authors conclude that ``high forecasting
power does not necessarily correspond to actionable trading signals''.

\begin{table}[ht]
\centering\footnotesize
\caption{Selected deep LOB models. F1 values are as reported by the authors on the data stated and
are not comparable across rows (Section~\ref{sec:comparable}).}
\label{tab:deep-models}
\setlength{\tabcolsep}{2pt}
\begin{tabularx}{\linewidth}{l*{4}{>{\raggedright\arraybackslash}X}}
\toprule
Model & Architecture & Data & Reported result & Independent re-test \\
\midrule
CNN~\citep{tsantekidis2017} & convolutions over time $\times$ levels & FI-2010 & F1 55--59\% & -- \\
LSTM~\citep{tsantekidis2017lstm} & 40-unit LSTM & FI-2010 & F1 61--66\% & -- \\
DeepLOB~\citep{zhang2019} & CNN + Inception + LSTM, $\sim$60k params & LSE, 5 stocks, 2017 & F1 61--70\% & re-trained on Nasdaq~\citep{briola2024}: $p_T$ 0.01--0.19 \\
TABL~\citep{tran2019tabl} & bilinear + temporal attention, $\sim$10k params & FI-2010 & F1 73--74\% & 70\% on FI-2010, 60\% on 2021 data~\citep{prata2024} \\
TransLOB~\citep{wallbridge2020} & dilated causal CNN + Transformer & FI-2010 & F1 81--92\% & 59\% on FI-2010~\citep{prata2024} \\
Axial-LOB~\citep{kisiel2022axial} & axial attention, $<$10k params & FI-2010 & F1 76--86\% & 73\% on FI-2010~\citep{prata2024} \\
TLOB~\citep{berti2025tlob} & dual spatial/temporal attention & FI-2010, Nasdaq, BTC & F1 to 93\% (FI-2010); 50--80\% (INTC) & -- \\
LiT~\citep{lit2025} & structured patches + Transformer & Binance, 2024 & F1 59--66\% & -- \\
OF-MATNet~\citep{ofmatnet} & multi-asset, multi-axis attention on OFI & 110 Nasdaq assets & $R^2$ gains in $>$90\% of cases & -- \\
\bottomrule
\end{tabularx}
\end{table}

\subsection{Representation versus architecture}

The benchmark evidence suggests that the choice and treatment of inputs can matter as much as the
network---a synthesis across studies, since none compares both systematically. Four decisions should be made, and reported, before capacity is increased.
\begin{description}[leftmargin=1.5em,style=nextline,itemsep=2pt]
\item[Representation] raw levels, order flow, volume on a price grid, or individual
  messages~\citep{lucchese2024,kolm2023,wu2021robust,linna2025lobert}.
\item[Normalisation] fixed statistics from the training set, rolling-window statistics, or a learned
  normalisation~\citep{passalis2020dain}; the statistics of LOB data drift over days.
\item[Target] label construction, threshold relative to the spread, and horizon~\citep{berti2025tlob}.
\item[Clock] event time, a fixed number of book updates, or calendar time; and how events are
  time-stamped~\citep{kolm_westray}.
\end{description}

\begin{quote}
\textbf{Evaluation rule (strongest-baseline comparison).} A complex architecture should be compared
not only against another complex architecture, but against the strongest reasonable low-capacity model
on the best available representation.
\end{quote}

\begin{table}[ht]
\centering\small
\caption{The comparisons a deep-model study should include. Each row should be reported with
out-of-sample log loss and out-of-sample net P\&L under the same execution model. The comparison
that matters is row~3 against row~1, not row~3 against row~2.}
\label{tab:baseline-template}
\setlength{\tabcolsep}{4pt}\footnotesize
\begin{tabular}{llll}
\toprule
Model & Input & Parameters & Typical inference latency \\
\midrule
Logistic & $\OFI^{(w)}_t$, $w = 1,\dots,5$\,s; $\QI_t$ & $\sim 10$ & ns \\
CNN--LSTM (DeepLOB-type) & 10-level snapshots & $\sim 10^5$ & $\mu$s--ms \\
Transformer & raw snapshots & $\sim 10^6$ & ms \\
Transformer & message stream $+$ OFI & $\sim 10^6$ & ms \\
\bottomrule
\end{tabular}
\end{table}

\subsection{Prediction targets}

Superficially similar targets define different statistical problems:
\[
  \begin{gathered}
  \P(\Delta m_{t,H} > 0 \mid X_t),\quad \E[\Delta m_{t,H} \mid X_t],\quad
  \P(\Delta m_{t,H} > \theta \mid X_t),\\
  \P(T_{\text{move}} \le H \mid X_t),\quad \P(\Fill \mid X_t),\quad \E[\Pi \mid X_t],
  \end{gathered}
\]
where $\theta > 0$ is a move threshold and $T_{\text{move}} = \inf\{u > 0 : |m_{t+u} - m_t| \ge
\delta\}$ is the time until the mid first moves by one tick.
\inwords{in order: the probability that the mid goes up; the average size of the move; the
probability that it goes up by more than $\theta$; the probability that it moves at all within $H$;
the probability that our resting order is filled within $H$; and our expected profit. These are six
different questions, and a model good at one need not be good at another.}
They should not be compared as if interchangeable. Benchmark papers should report the target,
horizon, sampling clock and intended execution interpretation.

\subsection{Fill-probability models}
\label{sec:deep-fill}

Deep learning has also been applied to the execution problem directly. \citet{arroyo2024} estimate
the full distribution of the time until a limit order is filled, using a convolutional-Transformer
encoder of recent market activity and a decoder that outputs a monotone survival curve. Cancellations
are treated as censored observations, and the models are trained and compared by the right-censored
log-likelihood, a proper scoring rule; the authors explain why the Brier score and concordance index
are unsuitable here, because their assumptions fail for limit orders. In their study, queue position
enters through the data design---hypothetical orders placed at the back of the queue---rather than as an
input. KANFormer~\citep{kanformer} adds the order's queue position and features of the agents
placing orders as inputs, on CAC~40 index futures, and reports a large improvement in both the
likelihood and discrimination measures over the earlier architecture, which it attributes to including
agent-level information and queue position. The study covers one futures contract, but it is
consistent with the claim of Section~\ref{sec:queue-position} that queue position is a first-class
variable.

\begin{modelcard}{Fill-probability (survival) models}
\cardrow{What you calibrate}{A survival curve $\P(\tau_F > u \mid X_t)$ for a resting order: for deep
models, the weights of an encoder and a decoder constrained to output a monotone curve.}
\cardrow{How}{Maximise the right-censored log-likelihood, treating cancelled orders as censored
observations; queue position enters as an input or through the design of hypothetical orders.}
\cardrow{What you do with it}{Choose between a limit and a market order, set how long to wait before
repricing, and weight expected profit by the probability of being filled.}
\cardrow{How to picture it}{A waiting-room clock for the order, which runs faster when the queue ahead
is short and the other side is aggressive.}
\end{modelcard}

\subsection{Generative and foundation models}
\label{sec:deep-gen}

\paragraph{Generators.} A growing literature generates synthetic order flow. \citet{nagy2023}
tokenise LOB messages and train an autoregressive model that generates message streams, processed
through a simulator; LOB-Bench~\citep{lobbench2025} standardises the evaluation of such generators.
Conditional generative adversarial networks act as ``world agents'' that produce realistic, reactive
order flow inside the ABIDES agent-based simulator~\citep{coletta2021,coletta2022,byrd2020abides}.
Diffusion models---which learn to reverse a gradual noising process---generate future volume
snapshots~\citep{diffvolume}, order flow conditioned on the book~\citep{berti2025trades}, and books
conditioned on a chosen future regime for counterfactual analysis~\citep{wang2026difflob}. M3
\citep{m3lob} trains an autoregressive model on a joint sequence of order events and book states for
800 Chinese stocks and about 32 billion events, reporting a scaling law across model sizes,
reproduction of stylised facts, and a square-root impact law when synthetic execution programmes are
injected. Neural extensions of queue-reactive models~\citep{deepqr} reproduce square-root impact and
cross-queue correlations on Bund futures.

\paragraph{How to judge a generator.} A generator is useful for research only if it preserves the
joint structure of orders, queues, executions and prices---not just marginal statistics such as the
average spread. The evaluation criteria in this literature are: marginal distributions (spreads,
depths, inter-event times); temporal dependence (clustering, long memory); cross-event dependence
(how cancellations follow trades); price dynamics (volatility, signature plots); liquidity dynamics
(queue depletion and refill); responsiveness to injected orders (market impact); and downstream
utility---whether a model trained on generated data performs well on real data~\citep{lobbench2025,jain2024sim}.
Consistent with Principle~1, this survey treats generators as tools for simulation and stress testing
and evaluates trading models only on observed data.

\begin{modelcard}{Generative models of order flow}
\cardrow{What you calibrate}{The weights of an autoregressive, adversarial or diffusion network that
produces messages or book states.}
\cardrow{How}{Next-token likelihood for autoregressive models, an adversarial game against a
discriminator for GANs, a denoising loss for diffusion models; then the generator is checked against
the criteria above.}
\cardrow{What you do with it}{Stress tests, counterfactual scenarios and training environments; not as
evidence that a trading model works.}
\cardrow{How to picture it}{A flight simulator for the order book: good for practice and rare
scenarios, not a substitute for flying.}
\end{modelcard}

\subsection{Foundation encoders and uncertainty}

A \emph{foundation model} is pretrained once on a large, generic corpus and then adapted to many
specific tasks. Two recent contributions from the same group illustrate a useful layering for LOB data.

\paragraph{LOBERT.} LOBERT is a pretrained, fine-tunable message-level encoder~\citep{linna2025lobert}.
Each message is one token (293 distinct message tokens), with its price, volume and time kept as
continuous values alongside; the context is 512 messages. It is pretrained by \emph{masked message
modelling}---hiding messages and asking the model to fill them in, as BERT does with words---on Nasdaq
ITCH data for AAPL, INTC, MSFT and FB from May to September 2015 (about 470 million messages, with
millisecond time stamps). The next-message model has about 1.1 million parameters and predicts the
complete next message correctly 27.8\% of the time, against 6.1\% for a state-space baseline of similar
size. On mid-price direction the comparison with DeepLOB depends on how many forecasts are kept. At
full coverage, 100 messages ahead, the two tie at a macro-F1 of 0.55. When each model forecasts only
when its confidence exceeds 0.9, LOBERT reaches 0.88 on the 10\% of cases it keeps, and DeepLOB 0.61 on
the 28\% it keeps---scores on different subsets, so not a like-for-like comparison. Inference is 53\%
slower than DeepLOB (about 282 predictions per second on one V100 GPU with batch size one, 47\% of
DeepLOB's throughput). The authors list as limitations that realistic long generated sequences are not
established, that the model may lose track of order size and location when price levels move, that it
does not model individual order identifiers, and that the data's millisecond resolution covers time
only partly.

\paragraph{Uncertainty layers.} UQ-LOB is an encoder-agnostic module
that attaches to a pretrained LOB encoder and, in the spirit of attentive neural processes,
conditions each forecast on a context set of recently completed windows whose outcomes are already
known~\citep{manoharan2026uqlob}; the paper evaluates it with one encoder, a deep temporal
attention bilinear network, on Kraken Level-2 and Level-3 data for seven cryptocurrency pairs.
Its regression variant outputs a Gaussian over the future tick
displacement and its classification variant a distribution over down, up and stationary; both
expose a scalar confidence. On 5.2 billion events across seven cryptocurrency assets it reports
near-nominal 68\% interval coverage.

For this survey the point is the change of target: from a direction to a direction \emph{with a
calibrated confidence}. That is precisely the input an execution policy needs, because the
decision to place, keep or cancel a passive order depends on how sure the model is, not only on
which way it leans (Sections~\ref{sec:tradeability} and~\ref{sec:passive}).

\paragraph{General time-series foundation models.} A parallel line of work pretrains forecasters on
very large collections of generic time series and applies them \emph{zero-shot}, without training on
the target series. Chronos~\citep{chronos2024} scales each series by its mean absolute value, bins the
values into 4096 tokens and trains a language-model architecture with a cross-entropy loss on about 84
billion observations. TimesFM~\citep{timesfm2024} is a 200-million-parameter decoder over patches of 32
time points, trained on about 100 billion points, most of them Wikipedia page views. Moirai
\citep{moirai2024} is a masked encoder that outputs a mixture distribution, trained on 27 billion
observations of which 0.1\% are economic or financial. Lag-Llama~\citep{lagllama2023} feeds a
decoder with lagged values and outputs a Student-$t$ distribution. None of these is trained or
evaluated on order-book data.

Finance-specific versions change the corpus. Kronos~\citep{kronos2025} is pretrained on over 12
billion candlestick bars (open, high, low, close, volume) from 45 exchanges and reports large gains
over general models in zero-shot price forecasting; its data are bars from five minutes to a day, not
messages. TradeFM~\citep{tradefm2026}, a 524-million-parameter preprint from J.P.\ Morgan AI Research,
tokenises individual trade-flow events (inter-arrival time, price depth, volume, side, add or cancel)
from more than 9{,}000 US equities and generates order flow inside a deterministic matching simulator.
Its rollouts reproduce heavy tails and volatility clustering, with a smaller distance to the observed
10-second return distribution than a compound Hawkes model (0.013 against 0.039), but the Hawkes model
matches the distribution of spreads better. The evidence for general models on financial prediction is
weak: on daily excess returns from 94 countries, Chronos and TimesFM used zero-shot have negative
out-of-sample $R^2$ ($-1.37\%$ and $-2.80\%$) and lose to gradient-boosted trees, and pretraining from
scratch on financial data closes much of the gap~\citep{rahimikia2025tsfm}. We found no fully reported
study of zero-shot general models on intraday or tick-level order-book data.

\begin{modelcard}{Foundation encoders and uncertainty layers}
\cardrow{What you calibrate}{An encoder pretrained on large volumes of message data; then a small task
head---and, for UQ-LOB, an uncertainty module---fine-tuned on the target task.}
\cardrow{How}{Self-supervised pretraining (for example, predicting masked or next messages), followed by
supervised fine-tuning; for the uncertainty layer, a context set of recently completed windows.}
\cardrow{What you do with it}{Start new tasks from a pretrained representation, forecast only when
confident (as in LOBERT's selective evaluation), and pass a calibrated confidence to an execution
policy that can threshold it. General time-series foundation models serve as baselines rather than
as LOB models.}
\cardrow{How to picture it}{A reader who has studied a library of order flow before being asked a
specific question, and who also says how sure the answer is.}
\end{modelcard}

% ================================================================
\section{Hybrid models}
\label{sec:hybrid}
% ================================================================

Hybrid models combine statistical microstructure structure with learned representations. The
motivation is that a purely learned model must rediscover from noisy data facts that are already known:
that events excite further events, that queues deplete and prices move when they do, that price impact
is roughly linear in order flow over short windows. Supplying this structure explicitly should make a
model more sample-efficient, more stable outside its training regime, and easier to interpret. These
are hypotheses; the published evidence for them in LOB data is still thin, and each hybrid should be
tested against its parts.

\subsection{What the statistical layer supplies}

The statistical models of the preceding sections provide low-dimensional, interpretable features that
can be computed in real time:
\begin{itemize}[leftmargin=1.5em,itemsep=1pt]
\item Hawkes intensities $\lambda_i(t)$ for each event type, updated in constant time per event;
\item best-level, multi-level and integrated OFI, and cross-asset OFI (Section~\ref{sec:ofi});
\item queue imbalance, the micro-price adjustment $m^{\text{micro}}_t - m_t$, and book-shape features
  (Section~\ref{sec:shape});
\item estimated depletion times $\tau^a, \tau^b$ and the up-move probability of the queue race
  (Section~\ref{sec:queue-models});
\item for execution models, the order's own queue position $Q^{\text{ahead}}_t$.
\end{itemize}

\subsection{Ways to combine the two}

\paragraph{Concatenation.} The simplest hybrid appends the statistical features to the network's
input, or projects them into the same embedding space as the learned features. It is easy to test
by ablation (below), and is the form most hybrids in the literature take.

\paragraph{Hawkes-biased attention.} Attention computes, for each event $n$, a weighted average over
earlier events $n'$, with weights proportional to the exponential of a score. A Hawkes kernel can be
built into that score:
\begin{equation}
  \text{score}(n, n') = \frac{\langle \mathbf{z}^{\text{qry}}_n, \mathbf{z}^{\text{key}}_{n'} \rangle}{\sqrt{\dim_{\text{emb}}}}
  + \ln \phi\big(t_n - t_{n'}\big), \qquad t_{n'} \le t_n ,
\end{equation}
where $\mathbf{z}^{\text{qry}}_n$ and $\mathbf{z}^{\text{key}}_{n'}$ are the query and key vectors of the
two events and $\dim_{\text{emb}}$ their dimension, as in Section~\ref{sec:deep}.
\inwords{because the softmax exponentiates the score, adding $\ln\phi$ \emph{multiplies} each
attention weight by the Hawkes kernel: recent events get more weight, and the weight decays with
elapsed time exactly as excitation does in a Hawkes process. With an exponential kernel
$\phi(u) = \alpha \mathrm{e}^{-\omega u}$ the added term is $\ln\alpha - \omega\,(t_n - t_{n'})$, a penalty linear
in elapsed time.} The network learns what to attend to; the kernel imposes how quickly influence
fades. Transformer Hawkes processes~\citep{zuo2020thp,zhang2020sahp} encode time differently, through
learned temporal encodings; the explicit kernel bias above is a design we propose for testing rather
than a published LOB result. A minimal implementation:
\begin{lstlisting}[language=Python]
def hawkes_biased_attention(E, t, Wq, Wk, Wv, log_alpha, omega):
    # E: one embedding row per event; t: event times (sorted)
    Q, K, V = E @ Wq, E @ Wk, E @ Wv
    scores = Q @ K.T / sqrt(Q.shape[1])
    dt = t[:, None] - t[None, :]              # elapsed time from event n' to event n
    scores = scores + (log_alpha - omega * dt)   # log of the exponential kernel
    scores[dt < 0] = -inf                     # causal: no attention to later events
    return softmax(scores, axis=1) @ V
\end{lstlisting}
The added cost over standard attention is one subtraction per pair of events.

\paragraph{Intensity-conditioned state-space models.} Selective state-space models such as
Mamba~\citep{gu2024mamba} let the step size of their recurrence depend on the input, so that the model
decides how quickly to forget. Making that step depend on the current event intensity or OFI would let
the model's memory shorten when the market is busy and lengthen when it is quiet---the same adaptation
that event-time sampling achieves by construction. We know of no published LOB model of this kind.

\paragraph{Residual gating.} A statistical model produces a baseline forecast $\widehat{Y}^{\text{stat}}_t$,
a network produces $\widehat{Y}^{\text{deep}}_t$, and a gate $\nu_t \in [0,1]$---learned, or driven by
volatility or intensity---mixes them:
\begin{equation}
  \widehat{Y}_t = \nu_t\, \widehat{Y}^{\text{stat}}_t + (1 - \nu_t)\, \widehat{Y}^{\text{deep}}_t .
\end{equation}
\inwords{when the gate is near 1 the forecast falls back on the simple, well-understood model; the
network contributes only where it has earned trust. Setting $\nu_t$ by volatility gives a natural safety
layer: in unusual conditions, defer to the model with fewer parameters.} A related form trains the
network on the residuals of the statistical model, so that it learns only what the simple model misses.

\subsection{Published hybrids}

Neural point processes for LOB events are the most developed hybrid. \citet{shi2022sdpnhp} compute a
separate intensity for each event type with parallel continuous-time LSTMs and add an interaction
between events and book state; on three US stocks it outperforms classical and neural
point-process baselines on likelihood and event prediction. \citet{lalor2025nhp} drive a twelve-event
LOB simulator with an LSTM-based neural Hawkes process and use it as the environment for a
reinforcement-learning market maker. The deep queue-reactive model of \citet{deepqr} keeps the
queue-reactive structure of \citet{huang2015} but parameterises the state-dependent event intensities
with neural networks. The uncertainty module UQ-LOB~\citep{manoharan2026uqlob} is hybrid in a different
sense: it attaches a probabilistic, attentive neural-process layer to a pretrained encoder.

\begin{modelcard}{Hybrid models}
\cardrow{What you calibrate}{The statistical component (for example Hawkes parameters or an OFI slope),
the network, and the way they are joined: a kernel bias in attention, or a gate $\nu_t$.}
\cardrow{How}{Fit the statistical part first, then train the network on its features or residuals; or
train end to end with the statistical part as a constrained layer. Always alongside the ablation
below.}
\cardrow{What you do with it}{Forecasts that need less data, behave more predictably outside the
training regime, and fall back on the simple model when the gate says so.}
\cardrow{How to picture it}{A trainee working under a checklist: the network adds judgement, the
statistical model supplies rules it should not have to rediscover.}
\end{modelcard}

\begin{example}[The right ablation]
Let $S_t$ be a Hawkes intensity computed from the same events already supplied to a
neural model $M_0(X_t)$. The experiment is $M_0(X_t)$ against $M_1(X_t, S_t)$ with identical
splits, seeds and economic evaluation, reporting $\Delta V$ with a confidence interval. We argue,
but do not test here, that $\Delta V \approx 0$ when $S_t$ is a deterministic function of $X_t$ and
the network sees enough context to compute it. A hybrid earns its place when $S_t$ brings
information the network cannot reconstruct, or when it lets a much smaller network match a larger
one.
\end{example}

% ================================================================
\section{Cross-asset information}
\label{sec:cross}
% ================================================================

The book of one instrument can carry information about related ones. A simple cross-impact form is
\begin{equation}
  m^{(\kappa)}_t - m^{(\kappa)}_{t-w} = \sum_{\kappa'} \beta_{\kappa \kappa'}\, \OFI^{(w,\kappa')}_t + \varepsilon^{(\kappa)}_t ,
\end{equation}
where the superscripts $(\kappa)$ and $(\kappa')$ index instruments, $m^{(\kappa)}_t$ is the mid of instrument $\kappa$,
$\OFI^{(w,\kappa')}_t$ is the order-flow imbalance of instrument $\kappa'$, $\beta_{\kappa \kappa'}$ is the cross-impact
of instrument $\kappa'$'s flow on instrument $\kappa$'s mid, and $\varepsilon^{(\kappa)}_t$ is a residual.
\inwords{the same straight-line fit as before, except that each instrument's price change may
respond to the buying pressure in \emph{every} instrument. The coefficients form a matrix; the
diagonal is each instrument's own impact and the off-diagonal entries are the cross effects.} The form covers
index constituents, ETF and constituents, futures and cash, correlated equities, related
fixed-income contracts and cross-venue cryptocurrency markets.

\begin{example}[Coefficients multiply]
For NQ,
\[
  m^{(\text{NQ})}_t - m^{(\text{NQ})}_{t-w} = \beta_{\text{NQ},\text{NQ}}\OFI^{(w,\text{NQ})}_t +
  \beta_{\text{NQ},\text{ES}}\OFI^{(w,\text{ES})}_t + \varepsilon^{(\text{NQ})}_t
\]
has two coefficients. Ten related instruments, each predicted from all ten OFIs at five lags, gives
$10 \times 10 \times 5 = 500$ coefficients. At a 5\% level about 25 are ``significant'' under the
null alone.
\end{example}

\paragraph{What the evidence says.} On the 100 largest S\&P~500 stocks, cross-asset order flow adds
nothing to the \emph{contemporaneous} fit once multi-level order flow is integrated, but lagged
cross-asset order flow does improve one-minute-ahead return forecasts, with an effect that decays
within minutes~\citep{contcucuringu2023}. Multi-asset attention models such as OF-MATNet attend along
the asset axis directly and choose which peers to use by Granger-causality tests~\citep{ofmatnet}. For
futures, the natural cross-asset structure is economic rather than statistical: index futures and
their constituents, related equity index futures (ES and NQ), and the Treasury futures curve.

\paragraph{Lead--lag between futures and cash.} The clearest cross-asset structure in futures is that
one market moves first. For the S\&P~500 and Nasdaq-100, most price discovery happens in the E-mini
futures rather than in the cash index or ETF~\citep{hasbrouck2003}. At high frequency the lead is
measured in milliseconds and is bounded by the physical distance between data centres. Between Chicago,
where CME matches, and the New Jersey equity venues, the response of SPY to E-mini price changes took
7.25--7.95\,ms in April 2010, fell after a straighter fibre route was built, and was estimated at 4.2--5.2\,ms
over microwave links, against a speed-of-light minimum of 3.93\,ms; about 20\% of SPY trading could be
identified as a direct response to E-mini moves~\citep{laughlin2014}. Measuring how often trades in two
markets occur at a given time offset, \citet{dobrev2017} find activity peaks at $\pm 5$\,ms between
10-year Treasury futures and cash Treasuries, with leadership running both ways, a clear lead of the
E-mini over SPY at $+5$\,ms, and a lead of the E-mini over cash Treasuries but not the reverse. The
order-flow link between SPY and the E-mini is ``positive, but short-lived'': beyond one second the
correlations have little economic significance~\citep{garrison2019}.

\paragraph{Estimating the lag.} Two markets trade at different, irregular times, so the usual
correlation of returns sampled on a common grid is biased towards zero (the Epps effect). The
Hayashi--Yoshida estimator avoids the grid by summing products of returns over every pair of
overlapping intervals; shifting one series by a candidate lag and choosing the lag that maximises the
estimated correlation gives a consistent lead--lag estimator~\citep{hoffmann2013leadlag}. On European data,
\citet{huth2014} find with this method that the CAC~40 future leads a large constituent stock by about
0.6 seconds on average, and that the leader's past alone predicts the direction of the lagger's next
mid-quote move about 60\% of the time---yet a naive strategy that crosses the spread on that forecast
does not make money.

\paragraph{Cross-asset Hawkes models.} Point-process models express lead--lag as cross-excitation:
events in one market raise the intensity of events in another. The two-asset model of
\citet{bacry2013} reproduces the Epps effect on Euro-Bund and Euro-Bobl futures, and lead--lag
can be read off its second-order properties; \citet{dafonseca2017} fit a multi-asset Hawkes model whose
diffusive limit links high- and low-frequency correlation and which captures lead--lag on Eurex
assets.

\paragraph{The Treasury curve.} For interest-rate futures the related instruments are the points of the
yield curve. On daily inter-dealer data from 1992 to 1999, order-flow imbalances explain up to 26\% of
day-to-day yield changes on days without major announcements, and most of that effect comes from a
factor common to all maturities: net buying across the whole curve says more about the level of yields
than buying in one maturity~\citep{brandt2004}. At one-minute resolution, \citet{tomas2022} compare
cross-impact models on 10-year Treasury note futures and E-mini futures jointly and find that
models built on the covariance of returns and order flow, such as a multivariate Kyle model, fit out of
sample far better than a model with only direct (own-instrument) impact. We found no published
intraday cross-impact study covering the full Treasury futures curve.

Cross-asset models should therefore report in-sample and out-of-sample improvement, incremental
economic value, added latency, parameter count and stability across regimes.

\begin{modelcard}{Cross-asset models}
\cardrow{What you calibrate}{The cross-impact matrix $\beta_{\kappa\kappa'}$ and the lags at which it
acts; for a cross-asset Hawkes model, the cross-excitation kernels between instruments.}
\cardrow{How}{Regression of each instrument's price change on the order flow of all (with shrinkage or
a factor structure to control the number of coefficients); lead--lag by maximising a Hayashi--Yoshida
correlation over candidate lags; Hawkes cross-kernels by maximum likelihood.}
\cardrow{What you do with it}{Use the leader's move to anticipate the lagger---mainly to protect resting
orders in the lagging market from being picked off---and hedge across related contracts.}
\cardrow{How to picture it}{Thunder after lightning: the futures move is the flash and the cash market
the sound, a few milliseconds later; the lag is set by distance.}
\end{modelcard}
